\documentclass{article}
\usepackage{amsmath}
\usepackage{amsfonts}
\usepackage{amssymb}
\usepackage{geometry}
\usepackage{algorithm}
\usepackage{algorithmic}
\usepackage{graphicx}
\usepackage{tikz}
\usetikzlibrary{shapes,arrows,positioning}

\geometry{margin=1in}

\title{System2Module Forward Pass: Memory-Enhanced Reasoning Architecture}
\author{Mathematical Analysis}
\date{\today}

\begin{document}

\maketitle

\section{Notation Conventions}

\textbf{Index Notation:} We use Einstein summation convention where repeated indices are summed over, except in Hadamard products where all indices must appear on both sides to avoid implicit summation.

\textbf{Hadamard Product ($\odot$):} Element-wise multiplication between tensors of identical dimensions. In index notation, all indices appear on both sides (e.g., $A_{ij} \odot B_{ij}$).

\textbf{Broadcasting:} When dimensions don't match for Hadamard product, we use broadcasting notation (e.g., $A_{ij} \alpha_i$ means $\alpha_i$ is repeated for each column $j$).

\textbf{Softmax Representation:} For any matrix $\mathbf{M}$, we define its softmax representation as $\mathbf{M}_S = \text{softmax}(\mathbf{M}, \text{dim}=-1)$.

\textbf{Sigmoid Representation:} For any matrix $\mathbf{M}$, we define its sigmoid representation as $\mathbf{M}_\sigma = \sigma(\mathbf{M})$ where $\sigma(x) = \frac{1}{1 + e^{-x}}$.

\textbf{Function Notation:} We use $\text{RMSNorm}(\cdot)$ for RMS normalization, $\text{softmax}(\cdot)$ for softmax, and $\sigma(\cdot)$ for sigmoid consistently throughout.

\textbf{Overbar Notation:} Variables with overbars (e.g., $\bar{\mathbf{g}}$, $\bar{\mathbf{u}}$) represent batch-averaged quantities: $\bar{\mathbf{x}} = \frac{1}{B}\sum_{b=1}^{B} \mathbf{x}[b, :]$.

\section{Overview}

The System2Module implements Kahneman's System 2 thinking through a sophisticated transformer architecture with persistent working memory. The forward pass combines input injection, multi-layer processing with memory enhancement, and reasoning history maintenance.

\section{Mathematical Formulation}

\subsection{Input Processing}

Given inputs:
\begin{align}
\mathbf{H} &\in \mathbb{R}^{B \times S \times D} \quad \text{(hidden states)} \\
\mathbf{I} &\in \mathbb{R}^{B \times S \times D} \quad \text{(input injection)} \\
\mathbf{R} &\in \mathbb{R}^{B \times P \times D} \quad \text{(reasoning history, optional)}
\end{align}

where $B$ is batch size, $S$ is sequence length, $D$ is hidden dimension, and $P$ is previous sequence length.

\subsection{Input Injection}

The first operation combines hidden states with input injection:

\begin{equation}
\mathbf{H}^{(0)} = \mathbf{H} + \mathbf{I} \in \mathbb{R}^{B \times S \times D}
\end{equation}

\textbf{Design Decision:} Simple addition is used rather than gating. This choice assumes that input injection should be additive rather than selective, allowing the model to learn appropriate combinations through subsequent layers.

\subsection{Layer-wise Processing}

For each layer $l \in \{1, 2, \ldots, L\}$ where $L$ is the number of effective layers:

\subsubsection{Working Memory Enhancement}

The working memory mechanism operates on the current state $\mathbf{x} = \mathbf{H}^{(l-1)}$:

\begin{enumerate}
\item \textbf{Memory Attention Computation:}
\begin{equation}
\mathbf{A} = \mathbf{x} \mathbf{M}^T \in \mathbb{R}^{B \times S \times M}
\end{equation}
where $\mathbf{M} \in \mathbb{R}^{M \times D}$ is the learnable memory buffer with $M$ memory slots.

\item \textbf{Attention Weight Normalization:}
\begin{equation}
\mathbf{W}_S = \text{softmax}(\mathbf{A}, \text{dim}=-1) \in \mathbb{R}^{B \times S \times M}
\end{equation}

\item \textbf{Memory Retrieval:}
\begin{equation}
\mathbf{R}_{\text{mem}} = \mathbf{W}_S \mathbf{M} \in \mathbb{R}^{B \times S \times D}
\end{equation}

\item \textbf{State Enhancement:}
\begin{equation}
\mathbf{x}_{\text{enhanced}} = \mathbf{x} + \mathbf{R}_{\text{mem}} \in \mathbb{R}^{B \times S \times D}
\end{equation}
\end{enumerate}

\textbf{Design Decision:} Memory retrieval uses attention weights rather than gating. This allows for soft, weighted combination of memory slots rather than element-wise gating. The addition operation assumes memory should be additive rather than replacing current information.

\subsubsection{Memory Retrieval Mechanisms}

The working memory supports three retrieval strategies, selected via the $\text{retrieval\_strategy}$ parameter:

\paragraph{Strategy 1: Simple Retrieval}

When $\text{retrieval\_strategy} = \text{"simple"}$, the system performs a straightforward attention-based memory retrieval that computes similarity scores between the current input and stored memory patterns.

\textbf{Parameter Definitions:}
\begin{itemize}
\item $\mathbf{M} \in \mathbb{R}^{M \times D}$: Memory matrix (learnable parameter)
\end{itemize}

\begin{enumerate}
\item \textbf{Attention Computation:}
\begin{align}
\mathbf{A}_{\text{ret}} &= \mathbf{x} \mathbf{M}^T \in \mathbb{R}^{B \times S \times M} \\
A_{\text{ret},bsi} &= \sum_{j=1}^{D} x_{bsj} M_{ij}
\end{align}

The attention computation calculates raw similarity scores between each input token $\mathbf{x}_{bs}$ and every memory slot $\mathbf{M}_i$. The dot product $A_{\text{ret},bsi}$ measures how similar the current input is to the $i$-th memory pattern. Higher values indicate greater similarity, suggesting that the corresponding memory slot contains relevant information for the current context.

\item \textbf{Attention Weights:}
\begin{align}
\mathbf{W}_{\text{ret},S} &= \text{softmax}(\mathbf{A}_{\text{ret}}, \text{dim}=-1) \in \mathbb{R}^{B \times S \times M} \\
W_{\text{ret},S,bsi} &= \frac{\exp(A_{\text{ret},bsi})}{\sum_{i'=1}^{M} \exp(A_{\text{ret},bsi'})}
\end{align}

The softmax normalization converts raw similarity scores into a probability distribution over memory slots. Each weight $W_{\text{ret},S,bsi}$ represents the probability that memory slot $i$ is relevant for input token $s$ in batch $b$. The softmax ensures that $\sum_{i=1}^{M} W_{\text{ret},S,bsi} = 1$, creating a proper attention mechanism where the model can attend to multiple memory slots with varying degrees of importance.

\item \textbf{Memory Retrieval:}
\begin{align}
\mathbf{R}_{\text{mem}} &= \mathbf{W}_{\text{ret},S} \mathbf{M} \in \mathbb{R}^{B \times S \times D} \\
R_{\text{mem},bsj} &= \sum_{i=1}^{M} W_{\text{ret},S,bsi} M_{ij}
\end{align}

The final retrieval step computes a weighted average of all memory slots, where each memory slot contributes proportionally to its attention weight. This creates a context-aware memory representation $\mathbf{R}_{\text{mem}}$ that combines relevant information from multiple memory patterns. The resulting vector $\mathbf{R}_{\text{mem},bs}$ represents what the model "remembers" about the current input, synthesized from its stored knowledge base.

\textbf{Asymptotic Complexity:} $O(B \cdot S \cdot D \cdot M)$
\begin{itemize}
\item Attention computation: $O(B \cdot S \cdot D \cdot M)$
\item Softmax computation: $O(B \cdot S \cdot M)$
\item Memory retrieval: $O(B \cdot S \cdot M \cdot D)$
\end{itemize}

\textbf{Memory Complexity:} $O(B \cdot S \cdot M + M \cdot D)$
\begin{itemize}
\item Attention matrix: $O(B \cdot S \cdot M)$
\item Memory matrix: $O(M \cdot D)$
\item Intermediate tensors: $O(B \cdot S \cdot D)$
\end{itemize}

\item \textbf{RMS Normalization:}
\begin{align}
\mathbf{R}_{\text{mem}} &= \text{RMSNorm}(\mathbf{R}_{\text{mem}}) \in \mathbb{R}^{B \times S \times D} \\
R_{\text{mem},bsj} &= \frac{R_{\text{mem},bsj}}{\sqrt{\frac{1}{D}\sum_{k=1}^{D} R_{\text{mem},bsk}^2 + \epsilon}} \cdot w_j
\end{align}

RMS normalization stabilizes the retrieved memory representations by scaling each dimension based on the root mean square of the vector. While not strictly necessary for the simple retrieval strategy, it provides consistency across all retrieval methods and ensures stable training behavior.
where $w_j$ is the learnable scale parameter and $\epsilon = 10^{-5}$ is a small constant for numerical stability.
\end{enumerate}

\paragraph{Strategy 2: Attention-Based Retrieval}

When $\text{retrieval\_strategy} = \text{"attention"}$, the system employs a sophisticated multi-head attention mechanism that allows the model to attend to different aspects of the memory simultaneously, providing richer memory interactions than simple dot-product attention.

\textbf{Parameter Definitions:}
\begin{itemize}
\item $\mathbf{W}_q, \mathbf{W}_k, \mathbf{W}_v \in \mathbb{R}^{D \times (D/H) \times H}$: Query, key, value projection matrices
\item $\mathbf{W}_{\text{out}} \in \mathbb{R}^{D \times D}$: Output projection matrix
\item $\mathbf{M} \in \mathbb{R}^{M \times D}$: Memory matrix (learnable parameter)
\end{itemize}

\begin{enumerate}
\item \textbf{Multi-Head Attention:}
\begin{align}
\mathbf{R}_{\text{mem}}, \mathbf{W}_{\text{ret},S} &= \text{MultiHeadAttention}(\mathbf{x}, \mathbf{M}, \mathbf{M}) \in \mathbb{R}^{B \times S \times D} \times \mathbb{R}^{B \times S \times M} \\
\text{where: } \mathbf{Q} &= \mathbf{x} \mathbf{W}_q, \quad \mathbf{K} = \mathbf{M} \mathbf{W}_k, \quad \mathbf{V} = \mathbf{M} \mathbf{W}_v \\
R_{\text{mem},bsj} &= \frac{1}{H}\sum_{h=1}^{H} \sum_{i=1}^{M} W_{\text{ret},S,bsi,h} V_{ij,h} \\
\text{where } W_{\text{ret},S,bsi,h} &= \text{softmax}(W_{\text{ret},bsi,h}) \\
\text{and } W_{\text{ret},bsi,h} &= \frac{\sum_{k=1}^{D} Q_{bsk,h} K_{ik,h}}{\sqrt{D/H}}
\end{align}

The multi-head attention mechanism transforms the input into three different representations: queries ($\mathbf{Q}$), keys ($\mathbf{K}$), and values ($\mathbf{V}$). Each head $h$ learns to attend to different types of relationships between the input and memory. The scaling factor $\sqrt{D/H}$ prevents the attention scores from becoming too large, maintaining stable gradients during training. The final output $\mathbf{R}_{\text{mem}}$ is a weighted combination of memory values, where the weights are determined by the learned attention patterns across all heads.

\item \textbf{RMS Normalization:}
\begin{align}
\mathbf{R}_{\text{mem}} &= \text{RMSNorm}(\mathbf{R}_{\text{mem}}) \in \mathbb{R}^{B \times S \times D} \\
R_{\text{mem},bsj} &= \frac{R_{\text{mem},bsj}}{\sqrt{\frac{1}{D}\sum_{k=1}^{D} R_{\text{mem},bsk}^2 + \epsilon}} \cdot w_j
\end{align}

RMS normalization stabilizes the retrieved memory representations by scaling each dimension based on the root mean square of the vector. This normalization technique, used in modern language models like LLaMA, provides better training stability compared to Layer Normalization while being computationally more efficient. The learnable scale parameter $w_j$ allows the model to adapt the normalization to the specific requirements of each dimension, where $w_j$ is the learnable scale parameter and $\epsilon = 10^{-5}$ is a small constant for numerical stability.

\textbf{Asymptotic Complexity:} $O(B \cdot S \cdot D^2 + B \cdot S \cdot M \cdot D)$
\begin{itemize}
\item Attention computation: $O(B \cdot S \cdot D^2 + B \cdot S \cdot M \cdot D)$
\item RMS normalization: $O(B \cdot S \cdot D)$
\end{itemize}

\textbf{Memory Complexity:} $O(B \cdot S \cdot M + M \cdot D + B \cdot S \cdot D)$
\begin{itemize}
\item Attention weights: $O(B \cdot S \cdot M)$
\item Memory matrix: $O(M \cdot D)$
\item Query/Key/Value projections: $O(B \cdot S \cdot D)$
\item RMS normalization parameters: $O(D)$
\end{itemize}
\end{enumerate}

\paragraph{Strategy 3: Sophisticated Retrieval}

When $\text{retrieval\_strategy} = \text{"sophisticated"}$, the system implements an advanced memory retrieval mechanism that incorporates temporal decay, importance weighting, and multi-head attention to create a more biologically plausible and contextually aware memory system.

\textbf{Parameter Definitions:}
\begin{itemize}
\item $\boldsymbol{\alpha}_{\text{ret}} \in \mathbb{R}^{M}$: Learnable memory decay factors for retrieval
\item $\mathbf{W}_q, \mathbf{W}_k, \mathbf{W}_v \in \mathbb{R}^{D \times (D/H) \times H}$: Query, key, value projection matrices
\item $\mathbf{W}_{\text{out}} \in \mathbb{R}^{D \times D}$: Output projection matrix
\item $\mathbf{W}_{\text{imp}} \in \mathbb{R}^{D \times 1}$: Learnable importance projection matrix
\item $\mathbf{M} \in \mathbb{R}^{M \times D}$: Memory matrix (learnable parameter)
\end{itemize}

\begin{enumerate}
\item \textbf{Memory Decay for Retrieval:}
\begin{align}
\mathbf{M}_{\text{decayed}} &= \mathbf{M} \odot (\boldsymbol{\alpha}_{\text{ret}} \mathbf{1}^T) \in \mathbb{R}^{M \times D} \\
M_{\text{decayed},ij} &= M_{ij} \alpha_{\text{ret},i} \quad \text{(broadcasting: $\alpha_{\text{ret},i}$ repeated for each column $j$)}
\end{align}

The memory decay mechanism simulates the natural forgetting process observed in human memory. Each memory slot $i$ has its own decay factor $\alpha_{\text{ret},i} \in [0,1]$, where values closer to 0 indicate stronger forgetting. This allows the model to learn which memories should be retained longer and which should decay more quickly, creating a more realistic memory system that can prioritize important information while allowing less relevant memories to fade.

\item \textbf{Multi-Head Attention with Decayed Memory:}
\begin{align}
\mathbf{R}_{\text{mem}}, \mathbf{W}_{\text{ret},S} &= \text{MultiHeadAttention}(\mathbf{x}, \mathbf{M}_{\text{decayed}}, \mathbf{M}_{\text{decayed}}) \in \mathbb{R}^{B \times S \times D} \times \mathbb{R}^{B \times S \times M} \\
\text{where: } \mathbf{Q} &= \mathbf{x} \mathbf{W}_q, \quad \mathbf{K} = \mathbf{M}_{\text{decayed}} \mathbf{W}_k, \quad \mathbf{V} = \mathbf{M}_{\text{decayed}} \mathbf{W}_v \\
R_{\text{mem},bsj} &= \sum_{h=1}^{H} \sum_{i=1}^{M} W_{\text{ret},S,bsi,h} V_{ij,h} \\
\text{where } W_{\text{ret},S,bsi,h} &= \text{softmax}(W_{\text{ret},bsi,h}) \\
\text{and } W_{\text{ret},bsi,h} &= \frac{\sum_{k=1}^{D} Q_{bsk,h} K_{ik,h}}{\sqrt{D/H}}
\end{align}

The attention mechanism operates on the decayed memory, ensuring that the model's attention is influenced by the temporal relevance of stored information. This creates a dynamic memory system where older, less relevant memories have reduced influence on the current computation, while recent and important memories maintain their full impact.

\item \textbf{RMS Normalization:}
\begin{align}
\mathbf{R}_{\text{mem}} &= \text{RMSNorm}(\mathbf{R}_{\text{mem}}) \in \mathbb{R}^{B \times S \times D} \\
R_{\text{mem},bsj} &= \frac{R_{\text{mem},bsj}}{\sqrt{\frac{1}{D}\sum_{k=1}^{D} R_{\text{mem},bsk}^2 + \epsilon}} \cdot w_j
\end{align}

The RMS normalization step ensures that the retrieved memory representations have consistent magnitude across different inputs and memory states, preventing the decay mechanism from causing numerical instabilities during training.
where $w_j$ is the learnable scale parameter and $\epsilon = 10^{-5}$ is a small constant for numerical stability.

\item \textbf{Importance Weighting:}
\begin{align}
\mathbf{I}_{\text{weights},\sigma} &= \sigma(\mathbf{x} \mathbf{W}_{\text{imp}}) \in \mathbb{R}^{B \times S \times 1} \\
I_{\text{weights},\sigma,bs1} &= \sigma\left(\sum_{j=1}^{D} x_{bsj} W_{\text{imp},j1}\right) \\
\mathbf{R}_{\text{mem}} &= \mathbf{R}_{\text{mem}} \odot \mathbf{I}_{\text{weights},\sigma} \in \mathbb{R}^{B \times S \times D} \\
R_{\text{mem},bsj} &= R_{\text{mem},bsj} I_{\text{weights},\sigma,bs1} \quad \text{(broadcasting: $I_{\text{weights},\sigma,bs1}$ repeated for each column $j$)}
\end{align}

The importance weighting mechanism provides a final layer of contextual filtering, allowing the model to dynamically adjust the strength of retrieved memory based on the current input's characteristics. The sigmoid activation ensures that importance weights are bounded between 0 and 1, where values closer to 1 indicate that the retrieved memory is highly relevant to the current context. This creates a sophisticated gating mechanism that can suppress irrelevant memories while amplifying important ones, similar to how human attention works in cognitive tasks.

\textbf{Asymptotic Complexity:} $O(B \cdot S \cdot D^2 + B \cdot S \cdot M \cdot D)$
\begin{itemize}
\item Memory decay: $O(M \cdot D)$
\item Attention computation: $O(B \cdot S \cdot D^2 + B \cdot S \cdot M \cdot D)$
\item RMS normalization: $O(B \cdot S \cdot D)$
\item Importance weighting: $O(B \cdot S \cdot D)$
\end{itemize}

\textbf{Memory Complexity:} $O(B \cdot S \cdot M + M \cdot D + B \cdot S \cdot D)$
\begin{itemize}
\item Attention weights: $O(B \cdot S \cdot M)$
\item Memory matrix: $O(M \cdot D)$
\item Decayed memory: $O(M \cdot D)$
\item Query/Key/Value projections: $O(B \cdot S \cdot D)$
\item Importance weights: $O(B \cdot S)$
\item RMS normalization parameters: $O(D)$
\end{itemize}
\end{enumerate}

\textbf{Design Decisions:} 
\begin{itemize}
\item \textbf{Strategy Selection:} Three levels of sophistication allow users to choose appropriate complexity for their task.
\item \textbf{Memory Decay:} Sophisticated strategy includes separate decay factors for retrieval vs update.
\item \textbf{Importance Weighting:} Sophisticated strategy weights retrieved memory based on current state importance.
\item \textbf{Multi-Head Attention:} Attention and sophisticated strategies use multi-head attention for richer memory interactions.
\item \textbf{RMS Normalization:} Used instead of LayerNorm for better efficiency and stability, following modern LLM practices (LLaMA, PaLM, etc.).
\end{itemize}

\subsubsection{Memory Update Mechanisms}

The working memory supports three update strategies, selected via the $\text{update\_strategy}$ parameter:

\paragraph{Strategy 1: Simple Update}

When $\text{update\_strategy} = \text{"simple"}$, the system implements a straightforward memory update mechanism that uses the first token of each sequence as a representative and applies context-dependent updates based on the current memory state.

\textbf{Parameter Definitions:}
\begin{itemize}
\item $\mathbf{W}_{\text{gate}} \in \mathbb{R}^{D \times M}$: Learnable gate projection matrix
\item $\mathbf{W}_{\text{update}} \in \mathbb{R}^{2D \times D}$: Learnable update projection matrix
\item $\mathbf{M} \in \mathbb{R}^{M \times D}$: Memory matrix (learnable parameter)
\end{itemize}

\begin{enumerate}
\item \textbf{Representative Selection:}
\begin{align}
\mathbf{r} &= \mathbf{x}[:, 0, :] \in \mathbb{R}^{B \times D} \\
r_{bj} &= x_{b0j} \quad \text{(first token of each sequence)}
\end{align}

The representative selection uses the first token of each sequence as a summary of the current input. This choice prioritizes computational efficiency over information content, relying on the model's ability to learn appropriate representations through the loss function. The first token often contains important contextual information in many language tasks, making it a reasonable heuristic for memory updates.

\item \textbf{Update Gate Computation:}
\begin{align}
\mathbf{g}_\sigma &= \sigma(\mathbf{r} \mathbf{W}_{\text{gate}}) \in \mathbb{R}^{B \times M} \\
g_{\sigma,bi} &= \sigma\left(\sum_{j=1}^{D} r_{bj} W_{\text{gate},ji}\right) \quad \text{where } \sigma(x) = \frac{1}{1 + e^{-x}}
\end{align}

The gate computation determines how much each memory slot should be updated based on the current representative. The sigmoid activation ensures that gate values are bounded between 0 and 1, where values closer to 1 indicate that the corresponding memory slot should be significantly updated. This creates a selective update mechanism that can preserve important memories while allowing others to be modified.

\item \textbf{Memory Update Preparation:}
\begin{align}
\bar{\mathbf{m}} &= \frac{1}{M}\sum_{i=1}^{M} \mathbf{M}[i, :] \in \mathbb{R}^{D} \\
\bar{m}_j &= \frac{1}{M}\sum_{i=1}^{M} M_{ij} \\
\mathbf{u} &= \mathbf{W}_{\text{update}} \begin{bmatrix} \mathbf{r} \\ \bar{\mathbf{m}} \end{bmatrix} \in \mathbb{R}^{B \times D} \\
u_{bj} &= \sum_{k=1}^{2D} W_{\text{update},kj} \begin{bmatrix} r_{bk} \\ \bar{m}_k \end{bmatrix}
\end{align}

The update preparation creates a context-dependent update signal by combining the current representative with the average memory state. The concatenation $[\mathbf{r}, \bar{\mathbf{m}}]$ allows the model to learn different update strategies based on both the current input and the existing memory context. This enables the system to adapt its memory updates based on what it already knows, creating a more sophisticated learning mechanism than simple input-based updates.

\item \textbf{Vectorized Memory Update:}
\begin{align}
\bar{\mathbf{g}}_\sigma &= \frac{1}{B}\sum_{b=1}^{B} \mathbf{g}_\sigma[b, :] \in \mathbb{R}^{M} \\
\bar{g}_{\sigma,i} &= \frac{1}{B}\sum_{b=1}^{B} g_{\sigma,bi} \\
\bar{\mathbf{u}} &= \frac{1}{B}\sum_{b=1}^{B} \mathbf{u}[b, :] \in \mathbb{R}^{D} \\
\bar{u}_j &= \frac{1}{B}\sum_{b=1}^{B} u_{bj} \\
\mathbf{M} &\leftarrow (1 - \bar{\mathbf{g}}_\sigma \mathbf{1}^T) \odot \mathbf{M} + (\bar{\mathbf{g}}_\sigma \mathbf{1}^T) \odot (\mathbf{1} \bar{\mathbf{u}}^T) \\
M_{ij} &\leftarrow (1 - \bar{g}_{\sigma,i}) M_{ij} + \bar{g}_{\sigma,i} \bar{u}_j \quad \text{(broadcasting: $\bar{g}_{\sigma,i}$ and $\bar{u}_j$ create $M \times D$ matrices)}
\end{align}

The final update step applies a weighted combination of the old memory and the new update signal. The formula $M_{ij} \leftarrow (1 - \bar{g}_{\sigma,i}) M_{ij} + \bar{g}_{\sigma,i} \bar{u}_j$ implements a gated update where $\bar{g}_{\sigma,i}$ controls the interpolation between preserving the old memory and incorporating the new information. This creates a smooth, continuous update process that can learn to balance memory stability with adaptability.

\textbf{Asymptotic Complexity:} $O(B \cdot D \cdot M + B \cdot D^2)$
\begin{itemize}
\item Gate computation: $O(B \cdot D \cdot M)$
\item Update preparation: $O(B \cdot D^2 + M \cdot D)$
\item Memory update: $O(M \cdot D)$
\end{itemize}

\textbf{Memory Complexity:} $O(B \cdot M + M \cdot D + B \cdot D)$
\begin{itemize}
\item Gate matrix: $O(B \cdot M)$
\item Memory matrix: $O(M \cdot D)$
\item Update signal: $O(B \cdot D)$
\item Gate projection matrix: $O(D \cdot M)$
\item Update projection matrix: $O(2D \cdot D)$
\end{itemize}
\end{enumerate}

\paragraph{Strategy 2: Attention-Based Update}

When $\text{update\_strategy} = \text{"attention"}$, the system employs a sophisticated attention mechanism to determine both what information to store and how much to update each memory slot, creating a more contextually aware memory update process.

\textbf{Parameter Definitions:}
\begin{itemize}
\item $\mathbf{W}_q, \mathbf{W}_k, \mathbf{W}_v \in \mathbb{R}^{D \times (D/H) \times H}$: Query, key, value projection matrices
\item $\mathbf{W}_{\text{out}} \in \mathbb{R}^{D \times D}$: Output projection matrix
\item $\mathbf{M} \in \mathbb{R}^{M \times D}$: Memory matrix (learnable parameter)
\end{itemize}

\begin{enumerate}
\item \textbf{Memory Attention:}
\begin{align}
\mathbf{A}_{\text{mem}}, \mathbf{W}_{\text{attn},S} &= \text{MultiHeadAttention}(\mathbf{x}, \mathbf{M}, \mathbf{M}) \in \mathbb{R}^{B \times S \times D} \times \mathbb{R}^{B \times S \times M} \\
\text{where: } \mathbf{Q} &= \mathbf{x} \mathbf{W}_q, \quad \mathbf{K} = \mathbf{M} \mathbf{W}_k, \quad \mathbf{V} = \mathbf{M} \mathbf{W}_v \\
A_{\text{mem},bsj} &= \sum_{h=1}^{H} \sum_{i=1}^{M} W_{\text{attn},S,bsi} V_{ij,h} \\
\text{where } W_{\text{attn},S,bsi} &= \text{softmax}(W_{\text{attn},bsi}) \\
\text{and } W_{\text{attn},bsi} &= \frac{\sum_{k=1}^{D} Q_{bsk,h} K_{ik,h}}{\sqrt{D/H}}
\end{align}

The memory attention mechanism computes how much each input token should attend to each memory slot, creating both an enhanced representation $\mathbf{A}_{\text{mem}}$ and attention weights $\mathbf{W}_{\text{attn},S}$. This allows the model to identify which memory slots are most relevant to the current input, providing a more nuanced understanding of memory relevance than simple similarity measures.

\item \textbf{Update Signal:}
\begin{align}
\mathbf{s} &= \frac{1}{S}\sum_{s=1}^{S} \mathbf{A}_{\text{mem}}[:, s, :] \in \mathbb{R}^{B \times D} \\
s_{bj} &= \frac{1}{S}\sum_{s=1}^{S} A_{\text{mem},bsj}
\end{align}

The update signal aggregates the attention-enhanced representations across all sequence positions, creating a single representative vector that captures the most important information from the current input. This averaging process ensures that the update signal reflects the overall context of the sequence rather than being dominated by individual tokens.

\item \textbf{Attention-Based Gates:}
\begin{align}
\mathbf{g}_{\text{attn},\sigma} &= \sigma\left(\frac{1}{S}\sum_{s=1}^{S} \mathbf{W}_{\text{attn},S}[:, s, :]\right) \in \mathbb{R}^{B \times M} \\
g_{\text{attn},\sigma,bi} &= \sigma\left(\frac{1}{S}\sum_{s=1}^{S} W_{\text{attn},S,bsi}\right)
\end{align}

The attention-based gates determine update strengths by averaging the attention weights across sequence positions. This creates a more sophisticated gating mechanism where the update strength for each memory slot is directly related to how much attention the current input paid to that slot. Memory slots that received more attention will be updated more strongly, creating a contextually appropriate update pattern.

\item \textbf{Vectorized Memory Update:}
\begin{align}
\bar{\mathbf{g}}_{\text{attn},\sigma} &= \frac{1}{B}\sum_{b=1}^{B} \mathbf{g}_{\text{attn},\sigma}[b, :] \in \mathbb{R}^{M} \\
\bar{g}_{\text{attn},\sigma,i} &= \frac{1}{B}\sum_{b=1}^{B} g_{\text{attn},\sigma,bi} \\
\bar{\mathbf{s}} &= \frac{1}{B}\sum_{b=1}^{B} \mathbf{s}[b, :] \in \mathbb{R}^{D} \\
\bar{s}_j &= \frac{1}{B}\sum_{b=1}^{B} s_{bj} \\
\mathbf{M} &\leftarrow (1 - \bar{\mathbf{g}}_{\text{attn},\sigma} \mathbf{1}^T) \odot \mathbf{M} + (\bar{\mathbf{g}}_{\text{attn},\sigma} \mathbf{1}^T) \odot (\mathbf{1} \bar{\mathbf{s}}^T) \\
M_{ij} &\leftarrow (1 - \bar{g}_{\text{attn},\sigma,i}) M_{ij} + \bar{g}_{\text{attn},\sigma,i} \bar{s}_j \quad \text{(broadcasting: $\bar{g}_{\text{attn},\sigma,i}$ and $\bar{s}_j$ create $M \times D$ matrices)}
\end{align}

\textbf{Asymptotic Complexity:} $O(B \cdot S \cdot D^2 + B \cdot S \cdot M \cdot D)$
\begin{itemize}
\item Attention computation: $O(B \cdot S \cdot D^2 + B \cdot S \cdot M \cdot D)$
\item Update signal: $O(B \cdot S \cdot D)$
\item Memory update: $O(M \cdot D)$
\end{itemize}

\textbf{Memory Complexity:} $O(B \cdot S \cdot M + M \cdot D + B \cdot S \cdot D)$
\begin{itemize}
\item Attention weights: $O(B \cdot S \cdot M)$
\item Memory matrix: $O(M \cdot D)$
\item Attention output: $O(B \cdot S \cdot D)$
\item Query/Key/Value projections: $O(B \cdot S \cdot D)$
\item Gate matrix: $O(B \cdot M)$
\end{itemize}
\end{enumerate}

\paragraph{Strategy 3: Sophisticated Update}

When $\text{update\_strategy} = \text{"sophisticated"}$, the system implements the most advanced memory update mechanism that combines temporal decay, importance-weighted representation, and adaptive learning rates to create a biologically inspired memory system with sophisticated forgetting and consolidation processes.

\textbf{Parameter Definitions:}
\begin{itemize}
\item $\boldsymbol{\alpha} \in \mathbb{R}^{M}$: Learnable memory decay factors
\item $\mathbf{W}_{\text{imp}} \in \mathbb{R}^{D \times 1}$: Learnable importance projection matrix
\item $\mathbf{W}_q, \mathbf{W}_k, \mathbf{W}_v \in \mathbb{R}^{D \times (D/H) \times H}$: Query, key, value projection matrices
\item $\mathbf{W}_{\text{out}} \in \mathbb{R}^{D \times D}$: Output projection matrix
\item $\mathbf{W}_{\text{gate}} \in \mathbb{R}^{D \times M}$: Learnable gate projection matrix
\item $\mathbf{M} \in \mathbb{R}^{M \times D}$: Memory matrix (learnable parameter)
\end{itemize}

\begin{enumerate}
\item \textbf{Memory Decay:}
\begin{align}
\mathbf{M} &\leftarrow \mathbf{M} \odot (\boldsymbol{\alpha} \mathbf{1}^T) \\
M_{ij} &\leftarrow M_{ij} \alpha_i \quad \text{(broadcasting: $\alpha_i$ repeated for each column $j$)}
\end{align}
where $\mathbf{1} \in \mathbb{R}^{D}$ is a vector of ones, making $\boldsymbol{\alpha} \mathbf{1}^T \in \mathbb{R}^{M \times D}$ for proper Hadamard product dimensions.

The memory decay step simulates the natural forgetting process that occurs in biological memory systems. Each memory slot has its own decay factor $\alpha_i \in [0,1]$, allowing the model to learn which memories should be retained longer and which should decay more quickly. This creates a dynamic memory system where less important information naturally fades over time, while crucial memories are preserved.

\item \textbf{Importance-Weighted Representative:}
\begin{align}
\mathbf{W}_{\text{imp},S} &= \text{softmax}(\sigma(\mathbf{x} \mathbf{W}_{\text{imp}}), \text{dim}=1) \in \mathbb{R}^{B \times S} \\
W_{\text{imp},S,bs} &= \frac{\exp(\sigma(\sum_{j=1}^{D} x_{bsj} W_{\text{imp},j1}))}{\sum_{s'=1}^{S} \exp(\sigma(\sum_{j=1}^{D} x_{bs'j} W_{\text{imp},j1}))} \\
\mathbf{r}_{\text{imp}} &= \sum_{s=1}^{S} \mathbf{W}_{\text{imp},S}[:, s] \odot \mathbf{x}[:, s, :] \in \mathbb{R}^{B \times D} \\
r_{\text{imp},bj} &= \sum_{s=1}^{S} W_{\text{imp},S,bs} x_{bsj} \quad \text{(broadcasting: $W_{\text{imp},S,bs}$ repeated for each column $j$)}
\end{align}

The importance-weighted representative mechanism replaces the simple "first token" approach with a learned, context-aware selection process. The model learns to identify which parts of the input sequence are most important for memory updates, creating a representative that reflects the overall significance of the current input rather than just its beginning. This allows the system to focus on the most relevant information when updating memory.

\textbf{Role of $\mathbf{W}_{\text{imp},S}$ and $\mathbf{r}_{\text{imp}}$:}
\begin{itemize}
\item \textbf{$\mathbf{W}_{\text{imp},S}$ (importance weights):} Softmax-normalized attention weights that determine how much each sequence position contributes to the representative. Each row sums to 1, creating a probability distribution over sequence positions.
\item \textbf{$\mathbf{r}_{\text{imp}}$ (importance-weighted representative):} A single representative vector per batch item that combines all sequence positions according to their learned importance. This replaces the simple "first token" approach with a learned, context-aware selection mechanism.
\end{itemize}

\item \textbf{Attention-Based Memory Interaction:}
\begin{align}
\mathbf{A}_{\text{soph}}, \mathbf{W}_{\text{soph},S} &= \text{MultiHeadAttention}(\mathbf{r}_{\text{imp}}, \mathbf{M}, \mathbf{M}) \in \mathbb{R}^{B \times 1 \times D} \times \mathbb{R}^{B \times 1 \times M} \\
\text{where: } \mathbf{Q} &= \mathbf{r}_{\text{imp}} \mathbf{W}_q, \quad \mathbf{K} = \mathbf{M} \mathbf{W}_k, \quad \mathbf{V} = \mathbf{M} \mathbf{W}_v \\
A_{\text{soph},b1j} &= \sum_{h=1}^{H} \sum_{i=1}^{M} W_{\text{soph},S,b1i} V_{ij,h} \\
\text{where } W_{\text{soph},S,b1i} &= \text{softmax}(W_{\text{soph},b1i}) \\
\text{and } W_{\text{soph},b1i} &= \frac{\sum_{k=1}^{D} Q_{b1k,h} K_{ik,h}}{\sqrt{D/H}}
\end{align}

The attention-based memory interaction computes how much the importance-weighted representative should attend to each memory slot. This creates a sophisticated interaction pattern where the model can identify which memories are most relevant to the current important information, providing a more nuanced understanding of memory relationships than simple similarity measures.

\item \textbf{Adaptive Learning Rates:}
\begin{align}
\mathbf{g}_{\text{base},\sigma} &= \sigma(\mathbf{r}_{\text{imp}} \mathbf{W}_{\text{gate}}) \in \mathbb{R}^{B \times M} \\
g_{\text{base},\sigma,bi} &= \sigma\left(\sum_{j=1}^{D} r_{\text{imp},bj} W_{\text{gate},ji}\right) \\
\mathbf{g}_{\text{final}} &= \mathbf{g}_{\text{base},\sigma} \odot \mathbf{W}_{\text{soph},S}[:, 0, :] \in \mathbb{R}^{B \times M} \\
g_{\text{final},bi} &= g_{\text{base},\sigma,bi} W_{\text{soph},S,b0i} \quad \text{(true Hadamard product: all indices $b,i$ on both sides)} \\
\boldsymbol{\lambda} &= \mathbf{g}_{\text{final}} \odot (\mathbf{1} (1 - \boldsymbol{\alpha})^T) \in \mathbb{R}^{B \times M} \\
\lambda_{bi} &= g_{\text{final},bi} (1 - \alpha_i) \quad \text{(broadcasting: $\alpha_i$ repeated for each batch $b$)}
\end{align}

The adaptive learning rates create a sophisticated update mechanism that combines three factors: (1) the base gate from the importance-weighted representative, (2) the attention weights from memory interaction, and (3) the decay factors. The final learning rate $\lambda_{bi}$ is higher for memories that are both relevant (high attention) and not already decayed (high $\alpha_i$), creating a contextually appropriate update pattern that balances memory stability with adaptability.

\item \textbf{Vectorized Sophisticated Update:}
\begin{align}
\bar{\boldsymbol{\lambda}} &= \frac{1}{B}\sum_{b=1}^{B} \boldsymbol{\lambda}[b, :] \in \mathbb{R}^{M} \\
\bar{\lambda}_i &= \frac{1}{B}\sum_{b=1}^{B} \lambda_{bi} \\
\bar{\mathbf{u}}_{\text{soph}} &= \frac{1}{B}\sum_{b=1}^{B} \mathbf{u}[b, :] \in \mathbb{R}^{D} \\
\bar{u}_{\text{soph},j} &= \frac{1}{B}\sum_{b=1}^{B} u_{bj} \\
\mathbf{M} &\leftarrow (1 - \bar{\boldsymbol{\lambda}} \mathbf{1}^T) \odot \mathbf{M} + (\bar{\boldsymbol{\lambda}} \mathbf{1}^T) \odot (\mathbf{1} \bar{\mathbf{u}}_{\text{soph}}^T) \\
M_{ij} &\leftarrow (1 - \bar{\lambda}_i) M_{ij} + \bar{\lambda}_i \bar{u}_{\text{soph},j} \quad \text{(broadcasting: $\bar{\lambda}_i$ and $\bar{u}_{\text{soph},j}$ create $M \times D$ matrices)}
\end{align}

The final update step applies the sophisticated learning rates to create a contextually aware memory update. The formula $M_{ij} \leftarrow (1 - \bar{\lambda}_i) M_{ij} + \bar{\lambda}_i \bar{u}_{\text{soph},j}$ implements an adaptive interpolation where the learning rate $\bar{\lambda}_i$ is determined by the combination of importance, attention, and decay factors. This creates a memory system that can learn to balance stability with adaptability based on the contextual relevance of information.

\textbf{Asymptotic Complexity:} $O(B \cdot S \cdot D + B \cdot D^2 + B \cdot M \cdot D)$
\begin{itemize}
\item Importance computation: $O(B \cdot S \cdot D)$
\item Attention computation: $O(B \cdot D^2 + B \cdot M \cdot D)$
\item Gate computation: $O(B \cdot D \cdot M)$
\item Memory update: $O(M \cdot D)$
\end{itemize}

\textbf{Memory Complexity:} $O(B \cdot S + M \cdot D + B \cdot D + B \cdot M)$
\begin{itemize}
\item Importance weights: $O(B \cdot S)$
\item Memory matrix: $O(M \cdot D)$
\item Attention weights: $O(B \cdot M)$
\item Query/Key/Value projections: $O(B \cdot D)$
\item Gate matrices: $O(B \cdot M)$
\item Decay factors: $O(M)$
\end{itemize}
\end{enumerate}

\textbf{Design Decisions:} 
\begin{itemize}
\item \textbf{Strategy Selection:} Three levels of sophistication allow users to choose appropriate complexity for their task.
\item \textbf{Vectorization:} All strategies use batch-averaged updates for computational efficiency.
\item \textbf{Representative Selection:} Simple uses first token, sophisticated uses importance-weighted average.
\item \textbf{Memory Decay:} Only sophisticated strategy includes forgetting mechanisms.
\end{itemize}

\subsubsection{Sophisticated Attention}

The attention mechanism processes the memory-enhanced state:

\begin{enumerate}
\item \textbf{Query, Key, Value Projections:}
\begin{align}
\mathbf{Q} &= \mathbf{W}_q \mathbf{x}_{\text{enhanced}} \in \mathbb{R}^{B \times S \times D} \\
\mathbf{K} &= \mathbf{W}_k \mathbf{x}_{\text{enhanced}} \in \mathbb{R}^{B \times S \times D} \\
\mathbf{V} &= \mathbf{W}_v \mathbf{x}_{\text{enhanced}} \in \mathbb{R}^{B \times S \times D}
\end{align}

\item \textbf{Multi-Head Reshaping:}
\begin{align}
\mathbf{Q} &\in \mathbb{R}^{B \times H \times S \times D_h} \\
\mathbf{K} &\in \mathbb{R}^{B \times H \times S \times D_h} \\
\mathbf{V} &\in \mathbb{R}^{B \times H \times S \times D_h}
\end{align}
where $H$ is the number of heads and $D_h = D/H$.

\item \textbf{Rotary Position Embedding:}
\begin{align}
\mathbf{Q}' &= \mathbf{Q} \odot \cos(\boldsymbol{\Theta}) + \text{rotate}(\mathbf{Q}) \odot \sin(\boldsymbol{\Theta}) \\
\mathbf{K}' &= \mathbf{K} \odot \cos(\boldsymbol{\Theta}) + \text{rotate}(\mathbf{K}) \odot \sin(\boldsymbol{\Theta})
\end{align}
where $\boldsymbol{\Theta} \in \mathbb{R}^{S \times D_h}$ contains position-dependent frequencies.

\item \textbf{Scaled Dot-Product Attention:}
\begin{equation}
\mathbf{A}_{\text{attn},S} = \text{softmax}\left(\frac{\mathbf{Q}' \mathbf{K}'^T}{\sqrt{D_h}}\right) \in \mathbb{R}^{B \times H \times S \times S}
\end{equation}

\item \textbf{Attention Output:}
\begin{equation}
\mathbf{O}_{\text{attn}} = \mathbf{A}_{\text{attn},S} \mathbf{V} \in \mathbb{R}^{B \times H \times S \times D_h}
\end{equation}

\item \textbf{Reshape and Project:}
\begin{equation}
\mathbf{O}_{\text{attn}} \in \mathbb{R}^{B \times S \times D} \rightarrow \mathbf{W}_{\text{out}} \mathbf{O}_{\text{attn}} \in \mathbb{R}^{B \times S \times D}
\end{equation}

\item \textbf{Reasoning History Attention (if available):}
\begin{equation}
\mathbf{O}_{\text{hist}}, \_ = \text{MultiHeadAttention}(\mathbf{O}_{\text{attn}}, \mathbf{R}, \mathbf{R}) \in \mathbb{R}^{B \times S \times D}
\end{equation}

\item \textbf{Final Attention Output:}
\begin{equation}
\mathbf{O}_{\text{final}} = \mathbf{O}_{\text{attn}} + \mathbf{O}_{\text{hist}} \in \mathbb{R}^{B \times S \times D}
\end{equation}
\end{enumerate}

\textbf{Design Decision:} Reasoning history is combined via addition rather than concatenation or gating. This maintains dimensional consistency while allowing the model to learn appropriate combinations through residual connections.

\subsubsection{Residual Connection and Normalization}

\begin{equation}
\mathbf{x}_{\text{attn}} = \mathbf{x}_{\text{enhanced}} + \mathbf{O}_{\text{final}} \in \mathbb{R}^{B \times S \times D}
\end{equation}

\begin{equation}
\mathbf{x}_{\text{norm1}} = \text{RMSNorm}(\mathbf{x}_{\text{attn}}) \in \mathbb{R}^{B \times S \times D}
\end{equation}

where RMSNorm is defined as:
\begin{equation}
\text{RMSNorm}(\mathbf{x}) = \mathbf{x} \odot \frac{1}{\sqrt{\frac{1}{D}\sum_{i=1}^{D} x_i^2 + \epsilon}}
\end{equation}

\subsubsection{MLP Processing (SwiGLU)}

\begin{enumerate}
\item \textbf{Gate and Up Projections:}
\begin{align}
\mathbf{G} &= \mathbf{W}_{\text{gate}} \mathbf{x}_{\text{norm1}} \in \mathbb{R}^{B \times S \times D_{\text{int}}} \\
\mathbf{U} &= \mathbf{W}_{\text{up}} \mathbf{x}_{\text{norm1}} \in \mathbb{R}^{B \times S \times D_{\text{int}}}
\end{align}

\item \textbf{Gated Activation:}
\begin{equation}
\mathbf{O}_{\text{mlp}} = \mathbf{W}_{\text{down}} (\text{SiLU}(\mathbf{G}) \odot \mathbf{U}) \in \mathbb{R}^{B \times S \times D}
\end{equation}

\item \textbf{Residual Connection and Normalization:}
\begin{align}
\mathbf{x}_{\text{mlp}} &= \mathbf{x}_{\text{norm1}} + \mathbf{O}_{\text{mlp}} \in \mathbb{R}^{B \times S \times D} \\
\mathbf{x}_{\text{final}} &= \text{RMSNorm}(\mathbf{x}_{\text{mlp}}) \in \mathbb{R}^{B \times S \times D}
\end{align}
\end{enumerate}

\textbf{Design Decision:} SwiGLU uses element-wise multiplication ($\odot$) rather than addition for gating. This allows for more sophisticated non-linear combinations and has shown superior performance in modern language models.

\subsection{Reasoning History Update}

After processing all layers, if in training mode:

\begin{equation}
\mathbf{R}_{\text{new}} = \begin{cases}
\mathbf{H}^{(L)} & \text{if } \mathbf{R} = \text{None} \\
\text{concat}(\mathbf{R}, \mathbf{H}^{(L)})[:, -P_{\max}:, :] & \text{otherwise}
\end{cases}
\end{equation}

where $P_{\max}$ is the maximum history length.

\textbf{Design Decision:} Concatenation is used for history accumulation rather than averaging or gating. This preserves the temporal sequence of reasoning steps, allowing the model to attend to specific historical positions.

\section{Algorithm Summary}

\begin{algorithm}
\caption{System2Module Forward Pass with Memory Update and Retrieval Strategies}
\begin{algorithmic}[1]
\STATE \textbf{Input:} $\mathbf{H} \in \mathbb{R}^{B \times S \times D}$, $\mathbf{I} \in \mathbb{R}^{B \times S \times D}$, $\mathbf{R} \in \mathbb{R}^{B \times P \times D}$, $\text{update\_strategy} \in \{\text{"simple"}, \text{"attention"}, \text{"sophisticated"}\}$, $\text{retrieval\_strategy} \in \{\text{"simple"}, \text{"attention"}, \text{"sophisticated"}\}$
\STATE $\mathbf{H}^{(0)} \leftarrow \mathbf{H} + \mathbf{I}$
\FOR{$l = 1$ to $L$}
    \STATE $\mathbf{x} \leftarrow \mathbf{H}^{(l-1)}$
    \STATE $\mathbf{x}_{\text{enhanced}} \leftarrow \text{WorkingMemory}(\mathbf{x}, \text{update\_strategy}, \text{retrieval\_strategy})$
    \STATE $\mathbf{x}_{\text{attn}} \leftarrow \mathbf{x}_{\text{enhanced}} + \text{Attention}(\mathbf{x}_{\text{enhanced}}, \mathbf{R})$
    \STATE $\mathbf{x}_{\text{norm1}} \leftarrow \text{RMSNorm}(\mathbf{x}_{\text{attn}})$
    \STATE $\mathbf{x}_{\text{mlp}} \leftarrow \mathbf{x}_{\text{norm1}} + \text{SwiGLU}(\mathbf{x}_{\text{norm1}})$
    \STATE $\mathbf{H}^{(l)} \leftarrow \text{RMSNorm}(\mathbf{x}_{\text{mlp}})$
\ENDFOR
\IF{training}
    \STATE $\mathbf{R} \leftarrow \text{UpdateHistory}(\mathbf{H}^{(L)})$
\ENDIF
\STATE \textbf{Return:} $\mathbf{H}^{(L)}$
\end{algorithmic}
\end{algorithm}

\begin{algorithm}
\caption{WorkingMemory with Vectorized Updates and Retrieval}
\begin{algorithmic}[1]
\STATE \textbf{Input:} $\mathbf{x} \in \mathbb{R}^{B \times S \times D}$, $\text{update\_strategy}$, $\text{retrieval\_strategy}$, $\text{update\_memory} = \text{True}$
\IF{$\text{retrieval\_strategy} = \text{"simple"}$}
    \STATE $\mathbf{A} \leftarrow \mathbf{x} \mathbf{M}^T$ \COMMENT{Memory attention}
    \STATE $\mathbf{W}_S \leftarrow \text{softmax}(\mathbf{A}, \text{dim}=-1)$ \COMMENT{Attention weights}
    \STATE $\mathbf{R}_{\text{mem}} \leftarrow \mathbf{W}_S \mathbf{M}$ \COMMENT{Memory retrieval}
\ELSIF{$\text{retrieval\_strategy} = \text{"attention"}$}
    \STATE $\mathbf{R}_{\text{mem}}, \mathbf{W}_S \leftarrow \text{MultiHeadAttention}(\mathbf{x}, \mathbf{M}, \mathbf{M})$ \COMMENT{Multi-head attention retrieval}
    \STATE $\mathbf{R}_{\text{mem}} \leftarrow \text{RMSNorm}(\mathbf{R}_{\text{mem}})$ \COMMENT{RMS normalization}
\ELSIF{$\text{retrieval\_strategy} = \text{"sophisticated"}$}
    \STATE $\mathbf{M}_{\text{decayed}} \leftarrow \mathbf{M} \odot (\boldsymbol{\alpha}_{\text{ret}} \mathbf{1}^T)$ \COMMENT{Memory decay for retrieval}
    \STATE $\mathbf{R}_{\text{mem}}, \mathbf{W}_S \leftarrow \text{MultiHeadAttention}(\mathbf{x}, \mathbf{M}_{\text{decayed}}, \mathbf{M}_{\text{decayed}})$ \COMMENT{Multi-head attention with decayed memory}
    \STATE $\mathbf{R}_{\text{mem}} \leftarrow \text{RMSNorm}(\mathbf{R}_{\text{mem}})$ \COMMENT{RMS normalization}
    \STATE $\mathbf{I}_{\text{weights},\sigma} \leftarrow \sigma(\mathbf{x} \mathbf{W}_{\text{imp}})$ \COMMENT{Importance weights}
    \STATE $\mathbf{R}_{\text{mem}} \leftarrow \mathbf{R}_{\text{mem}} \odot \mathbf{I}_{\text{weights},\sigma}$ \COMMENT{Importance weighting}
\ENDIF
\STATE $\mathbf{x}_{\text{enhanced}} \leftarrow \mathbf{x} + \mathbf{R}_{\text{mem}}$ \COMMENT{State enhancement}
\IF{update\_memory}
    \IF{update\_strategy = "simple"}
        \STATE $\mathbf{r} \leftarrow \mathbf{x}[:, 0, :]$ \COMMENT{First token representative}
        \STATE $\mathbf{g}_\sigma \leftarrow \sigma(\mathbf{W}_{\text{gate}} \mathbf{r})$ \COMMENT{Update gates}
        \STATE $\mathbf{u} \leftarrow \mathbf{W}_{\text{update}} [\mathbf{r}; \bar{\mathbf{m}}]$ \COMMENT{Update signal}
        \STATE $\mathbf{M} \leftarrow (1 - \bar{\mathbf{g}}_\sigma \mathbf{1}^T) \odot \mathbf{M} + (\bar{\mathbf{g}}_\sigma \mathbf{1}^T) \odot (\mathbf{1} \bar{\mathbf{u}}^T)$ \COMMENT{Vectorized update}
    \ELSIF{update\_strategy = "attention"}
        \STATE $(\mathbf{A}_{\text{mem}}, \mathbf{W}_{\text{attn},S}) \leftarrow \text{MultiHeadAttention}(\mathbf{x}, \mathbf{M}, \mathbf{M})$
        \STATE $\mathbf{s} \leftarrow \frac{1}{S}\sum_{s=1}^{S} \mathbf{A}_{\text{mem}}[:, s, :]$ \COMMENT{Update signal}
        \STATE $\mathbf{g}_{\text{attn},\sigma} \leftarrow \sigma(\frac{1}{S}\sum_{s=1}^{S} \mathbf{W}_{\text{attn},S}[:, s, :])$ \COMMENT{Attention gates}
        \STATE $\mathbf{M} \leftarrow (1 - \bar{\mathbf{g}}_{\text{attn},\sigma} \mathbf{1}^T) \odot \mathbf{M} + (\bar{\mathbf{g}}_{\text{attn},\sigma} \mathbf{1}^T) \odot (\mathbf{1} \bar{\mathbf{s}}^T)$ \COMMENT{Vectorized update}
    \ELSIF{update\_strategy = "sophisticated"}
        \STATE $\mathbf{M} \leftarrow \mathbf{M} \odot (\boldsymbol{\alpha} \mathbf{1}^T)$ \COMMENT{Memory decay (row-wise scaling)}
        \STATE $\mathbf{W}_{\text{imp},S} \leftarrow \text{softmax}(\sigma(\mathbf{W}_{\text{imp}} \mathbf{x}), \text{dim}=1)$ \COMMENT{Importance weights}
        \STATE $\mathbf{r}_{\text{imp}} \leftarrow \sum_{s=1}^{S} \mathbf{W}_{\text{imp},S}[:, s] \odot \mathbf{x}[:, s, :]$ \COMMENT{Importance representative (broadcasting: $W_{\text{imp},S,bs}$ repeated for each column)}
        \STATE $(\mathbf{A}_{\text{soph}}, \mathbf{W}_{\text{soph},S}) \leftarrow \text{MultiHeadAttention}(\mathbf{r}_{\text{imp}}, \mathbf{M}, \mathbf{M})$
        \STATE $\boldsymbol{\lambda} \leftarrow \sigma(\mathbf{W}_{\text{gate}} \mathbf{r}_{\text{imp}}) \odot \mathbf{W}_{\text{soph},S}[:, 0, :] \odot (\mathbf{1} (1 - \boldsymbol{\alpha})^T)$ \COMMENT{Adaptive learning rates}
        \STATE $\mathbf{M} \leftarrow (1 - \bar{\boldsymbol{\lambda}} \mathbf{1}^T) \odot \mathbf{M} + (\bar{\boldsymbol{\lambda}} \mathbf{1}^T) \odot (\mathbf{1} \bar{\mathbf{u}}^T)$ \COMMENT{Vectorized sophisticated update}
    \ENDIF
\ENDIF
\STATE \textbf{Return:} $\mathbf{x}_{\text{enhanced}}$
\end{algorithmic}
\end{algorithm}

\section{Design Analysis}

\subsection{Combination Strategies}

\begin{itemize}
\item \textbf{Addition vs. Concatenation:} Addition is used for residual connections and memory enhancement, maintaining dimensional consistency while allowing learned combinations. Concatenation is used for history accumulation to preserve temporal information.

\item \textbf{Gating vs. Attention:} Memory retrieval uses attention for soft, weighted combinations. MLP processing uses gating (SwiGLU) for non-linear feature interactions.

\item \textbf{Representative Selection:} First token selection prioritizes efficiency over information content, relying on the loss function to adapt the choice.

\end{itemize}

\subsection{Memory Architecture}

The working memory mechanism implements a dynamic memory system with three update strategies:

\subsubsection{Strategy Comparison}

\begin{table}[h]
\centering
\begin{tabular}{|l|c|c|c|}
\hline
\textbf{Feature} & \textbf{Simple} & \textbf{Attention} & \textbf{Sophisticated} \\
\hline
Representative Selection & First token & Attention-weighted & Importance-weighted \\
Memory Decay & No & No & Yes \\
Attention Mechanism & No & Yes & Yes \\
Adaptive Learning Rates & No & No & Yes \\
Computational Cost & Low & Medium & High \\
Memory Change Magnitude & High & Medium & Low \\
\hline
\end{tabular}
\caption{Comparison of memory update strategies}
\end{table}

\subsubsection{Vectorization Benefits}

All strategies use vectorized operations for computational efficiency:
\begin{itemize}
\item \textbf{Batch Averaging:} Update signals are averaged across batch dimension
\item \textbf{Element-wise Operations:} Memory updates use broadcasting for efficiency
\item \textbf{No Nested Loops:} Eliminated inefficient Python loops
\item \textbf{GPU Acceleration:} Operations can run in parallel on GPU
\end{itemize}

\subsubsection{Retrieval Strategy Comparison}

\begin{table}[h]
\centering
\begin{tabular}{|l|c|c|c|}
\hline
\textbf{Feature} & \textbf{Simple} & \textbf{Attention} & \textbf{Sophisticated} \\
\hline
Attention Mechanism & Basic & Multi-head & Multi-head \\
Memory Decay & No & No & Yes \\
Importance Weighting & No & No & Yes \\
RMS Normalization & No & Yes & Yes \\
Computational Cost & Low & Medium & High \\
Retrieval Quality & Basic & Rich & Adaptive \\
\hline
\end{tabular}
\caption{Comparison of memory retrieval strategies}
\end{table}

\subsubsection{Performance Characteristics}

\begin{itemize}
\item \textbf{Simple Strategy:} Fastest execution, largest memory changes, minimal parameters
\item \textbf{Attention Strategy:} Balanced performance, uses entire sequence information
\item \textbf{Sophisticated Strategy:} Most controlled updates, includes forgetting mechanisms
\end{itemize}

\subsubsection{Complexity Comparison}

\begin{table}[h]
\centering
\small
\hyphenpenalty=10000
\begin{tabular}{|p{2.2cm}|p{3.5cm}|p{3.5cm}|p{2.8cm}|}
\hline
\textbf{Strategy} & \textbf{Time Complexity} & \textbf{Memory Complexity} & \textbf{Dominant Factor} \\
\hline
Simple Retrieval & $O(BSDM)$ & $O(BSM + MD)$ & Memory slots $\times$ dims \\
\hline
Attention Retrieval & $O(BSD^2 + BSMD)$ & $O(BSM + MD + BSD)$ & Attention computation \\
\hline
Sophisticated Retrieval & $O(BSD^2 + BSMD)$ & $O(BSM + MD + BSD)$ & Attention + importance \\
\hline
Simple Update & $O(BDM + BD^2)$ & $O(BM + MD + BD)$ & Gate computation \\
\hline
Attention Update & $O(BSD^2 + BSMD)$ & $O(BSM + MD + BSD)$ & Attention computation \\
\hline
Sophisticated Update & $O(BSD + BD^2 + BMD)$ & $O(BS + MD + BD + BM)$ & Importance + attention \\
\hline
\end{tabular}
\caption{Computational and memory complexity comparison across all strategies}
\end{table}

This design provides a flexible foundation for persistent memory with configurable sophistication levels while maintaining computational efficiency through vectorization.

\section{Vectorization Improvements}

\subsection{Problem with Original Implementation}

The original memory update implementation used nested loops:
\begin{verbatim}
for b in range(batch_size):
    for m in range(memory_size):
        memory[m] = (1 - gate[b,m]) * memory[m] + gate[b,m] * update[b]
\end{verbatim}

This approach had several problems:
\begin{itemize}
\item \textbf{Poor Performance:} Loops are much slower than vectorized operations
\item \textbf{No GPU Utilization:} Cannot take advantage of parallel processing
\item \textbf{Memory Inefficient:} Multiple small operations instead of large vectorized ones
\item \textbf{Not Scalable:} Performance degrades significantly with larger batch sizes
\end{itemize}

\subsection{Vectorized Solution}

The improved implementation uses batch-averaged vectorized operations:

\begin{equation}
\mathbf{M} \leftarrow (1 - \bar{\mathbf{g}} \mathbf{1}^T) \odot \mathbf{M} + (\bar{\mathbf{g}} \mathbf{1}^T) \odot (\mathbf{1} \bar{\mathbf{u}}^T) \quad \text{(broadcasting creates proper Hadamard dimensions)}
\end{equation}

where $\bar{\mathbf{g}} = \frac{1}{B}\sum_{b=1}^{B} \mathbf{g}[b, :]$ and $\bar{\mathbf{u}} = \frac{1}{B}\sum_{b=1}^{B} \mathbf{u}[b, :]$.

\subsection{Performance Results}

\begin{table}[h]
\centering
\begin{tabular}{|l|c|c|c|c|}
\hline
\textbf{Configuration} & \textbf{Simple} & \textbf{Attention} & \textbf{Sophisticated} \\
\hline
Small (1×10×64) & 0.136 ms & 0.617 ms & 0.603 ms \\
Medium (4×20×64) & 0.329 ms & 1.275 ms & 1.477 ms \\
Large (8×50×128) & 0.534 ms & 3.481 ms & 2.999 ms \\
XLarge (16×100×256) & 0.597 ms & 11.352 ms & 6.190 ms \\
\hline
\end{tabular}
\caption{Memory update performance (milliseconds per update)}
\end{table}

\subsection{Key Benefits}

\begin{itemize}
\item \textbf{Linear Scaling:} Performance scales much better with batch size
\item \textbf{GPU Acceleration:} Operations can run in parallel on GPU
\item \textbf{Memory Efficiency:} Single large operations instead of many small ones
\item \textbf{Code Quality:} Cleaner, more maintainable vectorized code
\end{itemize}

\appendix

\section{MultiHeadAttention Two-Array Return}

This appendix explains why the \texttt{nn.MultiheadAttention} function returns two arrays and how they are used in the memory update strategies.

\subsection{MultiHeadAttention Output Structure}

The \texttt{MultiHeadAttention} function returns two distinct outputs:

\begin{equation}
\mathbf{A}_{\text{mem}}, \mathbf{W}_{\text{attn},S} = \text{MultiHeadAttention}(\mathbf{x}, \mathbf{M}, \mathbf{M}) \in \mathbb{R}^{B \times S \times D} \times \mathbb{R}^{B \times S \times M}
\end{equation}

\subsection{Output 1: Attention Output ($\mathbf{A}_{\text{mem}}$)}

\textbf{Purpose:} Contains the memory-enhanced representation of the input.

\textbf{Dimension:} $\mathbb{R}^{B \times S \times D}$

\textbf{What it represents:} For each token in the input sequence $\mathbf{x}$, this contains a new representation that has integrated relevant information from memory $\mathbf{M}$.

\textbf{Usage:} This is the primary output that gets used as the update signal in the next step of the memory update process.

\subsection{Output 2: Attention Weights ($\mathbf{W}_{\text{attn},S}$)}

\textbf{Purpose:} Contains the attention scores showing how much each input token attends to each memory slot.

\textbf{Dimension:} $\mathbb{R}^{B \times S \times M}$

\textbf{What it represents:} For each batch item and sequence position, there's a vector of $M$ attention scores (one per memory slot).

\textbf{Usage:} These weights are used to determine which memory slots to update and by how much.

\subsection{Mathematical Breakdown}

The \texttt{MultiHeadAttention(x, M, M)} operation works as follows:

\begin{itemize}
\item \textbf{Query:} $\mathbf{x}$ (current state) - "What am I looking for?"
\item \textbf{Key:} $\mathbf{M}$ (memory buffer) - "What's available in memory?"
\item \textbf{Value:} $\mathbf{M}$ (memory buffer) - "What information should I retrieve?"
\end{itemize}

\textbf{Attention Computation:}
\begin{equation}
\text{Attention}(Q, K, V) = \text{softmax}\left(\frac{QK^T}{\sqrt{d_k}}\right)V
\end{equation}

Where:
\begin{itemize}
\item $Q = \mathbf{x}$ (queries)
\item $K = \mathbf{M}$ (keys)
\item $V = \mathbf{M}$ (values)
\end{itemize}

\subsection{Why Both Arrays Are Needed}

\subsubsection{For Memory Enhancement:}
\begin{itemize}
\item $\mathbf{A}_{\text{mem}}$ provides the enhanced input representation
\item This gets averaged across sequence length to create the update signal
\end{itemize}

\subsubsection{For Memory Update:}
\begin{itemize}
\item $\mathbf{W}_{\text{attn},S}$ provides the attention weights
\item These weights determine which memory slots were most relevant
\item Used to compute update gates: $\mathbf{g}_{\text{attn},\sigma} = \sigma(\text{mean}(\mathbf{W}_{\text{attn},S}))$
\end{itemize}

\subsection{Usage in Memory Update Strategies}

\subsubsection{Attention-Based Strategy:}
\begin{enumerate}
\item \textbf{Step 1:} Compute attention between input and memory
   \begin{equation}
   \mathbf{A}_{\text{mem}}, \mathbf{W}_{\text{attn},S} = \text{MultiHeadAttention}(\mathbf{x}, \mathbf{M}, \mathbf{M})
   \end{equation}

\item \textbf{Step 2:} Use $\mathbf{A}_{\text{mem}}$ to create update signal
   \begin{equation}
   \mathbf{s} = \mathbf{A}_{\text{mem}}.\text{mean}(\text{dim}=1) \quad \text{(Average across sequence)}
   \end{equation}

\item \textbf{Step 3:} Use $\mathbf{W}_{\text{attn},S}$ to create update gates
   \begin{equation}
   \mathbf{g}_{\text{attn},\sigma} = \sigma(\mathbf{W}_{\text{attn},S}.\text{mean}(\text{dim}=1)) \quad \text{(Attention-based gates)}
   \end{equation}

\item \textbf{Step 4:} Apply vectorized update
   \begin{equation}
   \mathbf{M} \leftarrow (1 - \bar{\mathbf{g}}_{\text{attn},\sigma} \mathbf{1}^T) \odot \mathbf{M} + (\bar{\mathbf{g}}_{\text{attn},\sigma} \mathbf{1}^T) \odot (\mathbf{1} \bar{\mathbf{s}}^T) \quad \text{(broadcasting creates proper Hadamard dimensions)}
   \end{equation}
   where $\bar{\mathbf{g}}_{\text{attn},\sigma} = \frac{1}{B}\sum_{b=1}^{B} \mathbf{g}_{\text{attn},\sigma}[b, :]$ and $\bar{\mathbf{s}} = \frac{1}{B}\sum_{b=1}^{B} \mathbf{s}[b, :]$
\end{enumerate}

\subsubsection{Sophisticated Strategy:}
The sophisticated strategy uses a similar pattern but with a single representative as the query:

\begin{enumerate}
\item \textbf{Step 1:} Compute attention between representative and memory
   \begin{equation}
   \mathbf{A}_{\text{soph}}, \mathbf{W}_{\text{soph},S} = \text{MultiHeadAttention}(\mathbf{r}_{\text{imp}}, \mathbf{M}, \mathbf{M})
   \end{equation}

\item \textbf{Step 2:} Use both outputs for sophisticated memory updates with decay and importance weighting
\end{enumerate}

\subsection{Key Insight}

The two outputs serve complementary roles:
\begin{itemize}
\item \textbf{$\mathbf{A}_{\text{mem}}$:} "What information was retrieved from memory?"
\item \textbf{$\mathbf{W}_{\text{attn},S}$:} "How much attention was paid to each memory slot?"
\end{itemize}

This dual output allows the attention-based strategy to both enhance the current state with memory information AND use the attention patterns to intelligently update the memory buffer itself.

\section{Architectural Design vs. Learned Behavior}

\subsection{The Gap Between Design Intent and Training Reality}

A critical consideration in the design of memory update mechanisms is the distinction between \textbf{what the architecture enables} and \textbf{what the model actually learns}. While the simple memory update strategy is designed to enable context-dependent learning through the formulation:

\begin{equation}
\mathbf{u} = \mathbf{W}_{\text{update}} \begin{bmatrix} \mathbf{r} \\ \bar{\mathbf{m}} \end{bmatrix}
\end{equation}

the actual behavior depends entirely on the learned weights $\mathbf{W}_{\text{update}} \in \mathbb{R}^{D \times 2D}$.

\subsection{The Hot Stove Analogy}

Consider the analogy of learning about fire safety through touching a hot stove:

\textbf{Baby's Experience:}
\begin{itemize}
\item Limited existing knowledge about heat, pain, and safety
\item Simple, direct association: ``hot stove → pain''
\item Memory update: Basic safety rule formation
\end{itemize}

\textbf{Teenager's Experience:}
\begin{itemize}
\item Rich existing knowledge about chemistry, physics, and safety protocols
\item Complex integration: ``hot stove → thermal energy transfer → cellular damage → pain response''
\item Memory update: Sophisticated understanding of heat transfer mechanisms
\end{itemize}

Both receive the same sensory input (touching the hot stove), but their memory integration strategies differ dramatically based on their existing knowledge base $\bar{\mathbf{m}}$.

\subsection{Theoretical vs. Practical Learning}

\textbf{Theoretical Capability:}
The architecture enables context-dependent updates through the concatenation $[\mathbf{r}, \bar{\mathbf{m}}]$, allowing the model to learn different integration strategies for different knowledge states.

\textbf{Practical Reality:}
The learned weights $\mathbf{W}_{\text{update}}$ are determined by:
\begin{itemize}
\item \textbf{Loss function design:} What behavior is rewarded during training
\item \textbf{Training data distribution:} What patterns the model observes
\item \textbf{Optimization success:} Whether the model finds optimal solutions
\end{itemize}

\subsection{Potential Learning Outcomes}

The model might learn:
\begin{itemize}
\item $\mathbf{W}_{\text{update}} \approx [\mathbf{W}_r, \mathbf{0}]$: Ignores memory context entirely
\item $\mathbf{W}_{\text{update}} \approx [\mathbf{0}, \mathbf{W}_m]$: Ignores new information entirely
\item $\mathbf{W}_{\text{update}} \approx [\mathbf{W}_r, \mathbf{W}_m]$: Uses both components (intended behavior)
\end{itemize}

\subsection{Empirical Validation Requirements}

To verify that context-dependent learning actually occurs, empirical analysis must examine:
\begin{enumerate}
\item \textbf{Weight analysis:} Whether $\mathbf{W}_{\text{update}}$ utilizes both input components
\item \textbf{Behavioral testing:} Whether updates vary with memory state
\item \textbf{Performance comparison:} Whether context-dependent updates improve task performance
\item \textbf{Ablation studies:} Whether removing memory context degrades performance
\end{enumerate}

\subsection{Design Implications}

This analysis highlights the importance of:
\begin{itemize}
\item \textbf{Careful loss function design} to encourage context-dependent behavior
\item \textbf{Diverse training data} that demonstrates context-dependent patterns
\item \textbf{Empirical validation} to verify that intended behaviors emerge
\item \textbf{Architectural constraints} that make desired behaviors more likely to be learned
\end{itemize}

The memory update mechanism provides the \textit{possibility} of sophisticated, context-dependent learning, but successful training is required to realize this potential.

\section{Appendix C: Relationships Between Retrieval Strategies}

This appendix analyzes the mathematical and conceptual relationships between the three memory retrieval strategies, demonstrating how they form a hierarchical progression from simple to sophisticated mechanisms.

\subsection{Strategy 2 as a Generalization of Strategy 1}

\subsubsection{Mathematical Relationship}

Strategy 1 (Simple Retrieval) can be viewed as a special case of Strategy 2 (Attention-Based Retrieval) when specific constraints are applied to the learned parameters.

\textbf{Strategy 1 Formulation:}
\begin{align}
\mathbf{A}_{\text{ret}} &= \mathbf{x} \mathbf{M}^T \\
\mathbf{W}_{\text{ret},S} &= \text{softmax}(\mathbf{A}_{\text{ret}}) \\
\mathbf{R}_{\text{mem}} &= \mathbf{W}_{\text{ret},S} \mathbf{M}
\end{align}

\textbf{Strategy 2 Formulation (Single Head):}
\begin{align}
\mathbf{Q} &= \mathbf{x} \mathbf{W}_q, \quad \mathbf{K} = \mathbf{M} \mathbf{W}_k, \quad \mathbf{V} = \mathbf{M} \mathbf{W}_v \\
\mathbf{A}_{\text{attn}} &= \frac{\mathbf{Q} \mathbf{K}^T}{\sqrt{D}} \\
\mathbf{W}_{\text{ret},S} &= \text{softmax}(\mathbf{A}_{\text{attn}}) \\
\mathbf{R}_{\text{mem}} &= \mathbf{W}_{\text{ret},S} \mathbf{V}
\end{align}

\textbf{Equivalence Conditions:}

Strategy 2 reduces to Strategy 1 when:
\begin{itemize}
\item \textbf{Identity Projections:} $\mathbf{W}_q = \mathbf{W}_k = \mathbf{W}_v = \mathbf{I}$
\item \textbf{Single Head:} $H = 1$
\item \textbf{Unit Scaling:} $\sqrt{D} = 1$ (or absorbed into learning)
\end{itemize}

Under these conditions:
\begin{equation}
\mathbf{Q} = \mathbf{x}, \quad \mathbf{K} = \mathbf{V} = \mathbf{M}
\end{equation}

This gives:
\begin{equation}
\mathbf{A}_{\text{attn}} = \mathbf{x} \mathbf{M}^T = \mathbf{A}_{\text{ret}}
\end{equation}

Therefore, Strategy 1 is Strategy 2 with \textbf{constrained projections}.

\subsubsection{Conceptual Interpretation}

\begin{itemize}
\item \textbf{Strategy 1:} "What is the direct similarity between input and memory?"
\item \textbf{Strategy 2:} "What is the learned similarity between transformed input and transformed memory?"
\end{itemize}

Strategy 2 generalizes Strategy 1 by allowing the model to learn \textbf{task-specific similarity functions} through the projection matrices $\mathbf{W}_q, \mathbf{W}_k, \mathbf{W}_v$.

\subsection{Strategy 3 as an Extension of Strategy 2}

\subsubsection{Mathematical Relationship}

Strategy 3 (Sophisticated Retrieval) extends Strategy 2 by adding three key components: memory decay, importance weighting, and enhanced memory interaction.

\textbf{Strategy 2 Base:}
\begin{align}
\mathbf{R}_{\text{mem}}, \mathbf{W}_{\text{ret},S} &= \text{MultiHeadAttention}(\mathbf{x}, \mathbf{M}, \mathbf{M}) \\
\mathbf{R}_{\text{mem}} &= \text{RMSNorm}(\mathbf{R}_{\text{mem}})
\end{align}

\textbf{Strategy 3 Extensions:}
\begin{enumerate}
\item \textbf{Memory Decay:}
\begin{equation}
\mathbf{M}_{\text{decayed}} = \mathbf{M} \odot (\boldsymbol{\alpha}_{\text{ret}} \mathbf{1}^T)
\end{equation}

\item \textbf{Decayed Memory Attention:}
\begin{equation}
\mathbf{R}_{\text{mem}}, \mathbf{W}_{\text{ret},S} = \text{MultiHeadAttention}(\mathbf{x}, \mathbf{M}_{\text{decayed}}, \mathbf{M}_{\text{decayed}})
\end{equation}

\item \textbf{Importance Weighting:}
\begin{align}
\mathbf{I}_{\text{weights},\sigma} &= \sigma(\mathbf{x} \mathbf{W}_{\text{imp}}) \\
\mathbf{R}_{\text{mem}} &= \mathbf{R}_{\text{mem}} \odot \mathbf{I}_{\text{weights},\sigma}
\end{align}
\end{enumerate}

\textbf{Reduction Conditions:}

Strategy 3 reduces to Strategy 2 when:
\begin{itemize}
\item \textbf{No Decay:} $\boldsymbol{\alpha}_{\text{ret}} = \mathbf{1}$ (all ones vector)
\item \textbf{No Importance Weighting:} $\mathbf{I}_{\text{weights},\sigma} = \mathbf{1}$ (all ones)
\end{itemize}

Under these conditions:
\begin{align}
\mathbf{M}_{\text{decayed}} &= \mathbf{M} \\
\mathbf{R}_{\text{mem}} &= \mathbf{R}_{\text{mem}} \odot \mathbf{1} = \mathbf{R}_{\text{mem}}
\end{align}

\subsubsection{Conceptual Interpretation}

Strategy 3 adds \textbf{biological realism} and \textbf{contextual awareness}:

\begin{itemize}
\item \textbf{Memory Decay:} Simulates natural forgetting processes where older or less important memories fade over time
\item \textbf{Importance Weighting:} Provides context-dependent gating that can suppress irrelevant memories based on current input characteristics
\item \textbf{Enhanced Attention:} Operates on temporally-aware memory representations
\end{itemize}

\subsection{Hierarchical Progression}

The three strategies form a \textbf{hierarchical progression} of increasing sophistication:

\begin{table}[h]
\centering
\begin{tabular}{|l|c|c|c|}
\hline
\textbf{Feature} & \textbf{Strategy 1} & \textbf{Strategy 2} & \textbf{Strategy 3} \\
\hline
Similarity Function & Fixed (dot product) & Learned (attention) & Learned (attention) \\
Memory Processing & Direct & Direct & Decayed \\
Output Processing & Direct & RMS normalized & RMS + importance \\
Temporal Awareness & No & No & Yes \\
Context Filtering & No & No & Yes \\
\hline
\end{tabular}
\caption{Hierarchical progression of retrieval strategies}
\end{table}

\subsection{Unified Mathematical Framework}

All three strategies can be expressed in a unified framework:

\begin{align}
\mathbf{M}' &= \mathbf{M} \odot (\boldsymbol{\alpha} \mathbf{1}^T) \quad &\text{(Memory decay)} \\
\mathbf{R}_{\text{mem}} &= \text{Attention}(\mathbf{x}, \mathbf{M}', \mathbf{M}') \quad &\text{(Similarity computation)} \\
\mathbf{R}_{\text{mem}} &= \text{Normalize}(\mathbf{R}_{\text{mem}}) \quad &\text{(Stabilization)} \\
\mathbf{R}_{\text{mem}} &= \mathbf{R}_{\text{mem}} \odot \mathbf{I}_{\text{weights}} \quad &\text{(Importance filtering)}
\end{align}

Where:
\begin{itemize}
\item \textbf{Strategy 1:} $\boldsymbol{\alpha} = \mathbf{1}$, Attention = dot product, Normalize = identity, $\mathbf{I}_{\text{weights}} = \mathbf{1}$
\item \textbf{Strategy 2:} $\boldsymbol{\alpha} = \mathbf{1}$, Attention = multi-head, Normalize = RMSNorm, $\mathbf{I}_{\text{weights}} = \mathbf{1}$
\item \textbf{Strategy 3:} $\boldsymbol{\alpha}$ = learned, Attention = multi-head, Normalize = RMSNorm, $\mathbf{I}_{\text{weights}}$ = learned
\end{itemize}

\subsection{Design Implications}

This hierarchical relationship has important implications:

\begin{enumerate}
\item \textbf{Progressive Complexity:} Each strategy builds upon the previous one, allowing for gradual sophistication
\item \textbf{Backward Compatibility:} More sophisticated strategies can degrade gracefully to simpler ones
\item \textbf{Ablation Studies:} The relationships enable systematic analysis of each component's contribution
\item \textbf{Adaptive Selection:} Systems can dynamically choose appropriate complexity based on task requirements
\end{enumerate}

The unified framework demonstrates that memory retrieval sophistication lies on a continuum, where each strategy represents a different point balancing computational cost, biological plausibility, and representational power.

\section{Appendix D: Relationships Between Memory Update Strategies}

This appendix analyzes the mathematical and conceptual relationships between the three memory update strategies, demonstrating how they form a hierarchical progression from simple to sophisticated memory storage mechanisms.

\subsection{Strategy 2 as a Generalization of Strategy 1}

\subsubsection{Mathematical Relationship}

Strategy 1 (Simple Update) can be viewed as a special case of Strategy 2 (Attention-Based Update) when specific constraints are applied to the attention mechanism and representative selection.

\textbf{Strategy 1 Formulation:}
\begin{align}
\mathbf{r} &= \mathbf{x}[:, 0, :] \quad \text{(First token representative)} \\
\mathbf{g}_\sigma &= \sigma(\mathbf{r} \mathbf{W}_{\text{gate}}) \\
\bar{\mathbf{m}} &= \frac{1}{M}\sum_{i=1}^{M} \mathbf{M}[i, :] \\
\mathbf{u} &= \mathbf{W}_{\text{update}} \begin{bmatrix} \mathbf{r} \\ \bar{\mathbf{m}} \end{bmatrix} \\
\mathbf{M} &\leftarrow (1 - \bar{\mathbf{g}}_\sigma \mathbf{1}^T) \odot \mathbf{M} + (\bar{\mathbf{g}}_\sigma \mathbf{1}^T) \odot (\mathbf{1} \bar{\mathbf{u}}^T)
\end{align}

\textbf{Strategy 2 Formulation:}
\begin{align}
\mathbf{A}_{\text{mem}}, \mathbf{W}_{\text{attn},S} &= \text{MultiHeadAttention}(\mathbf{x}, \mathbf{M}, \mathbf{M}) \\
\mathbf{s} &= \frac{1}{S}\sum_{s=1}^{S} \mathbf{A}_{\text{mem}}[:, s, :] \quad \text{(Attention-based representative)} \\
\mathbf{g}_{\text{attn},\sigma} &= \sigma\left(\frac{1}{S}\sum_{s=1}^{S} \mathbf{W}_{\text{attn},S}[:, s, :]\right) \\
\mathbf{M} &\leftarrow (1 - \bar{\mathbf{g}}_{\text{attn},\sigma} \mathbf{1}^T) \odot \mathbf{M} + (\bar{\mathbf{g}}_{\text{attn},\sigma} \mathbf{1}^T) \odot (\mathbf{1} \bar{\mathbf{s}}^T)
\end{align}

\textbf{Equivalence Conditions:}

Strategy 2 reduces to Strategy 1 when:
\begin{itemize}
\item \textbf{Identity Attention:} $\mathbf{W}_q = \mathbf{W}_k = \mathbf{W}_v = \mathbf{I}$ and $H = 1$
\item \textbf{First Token Focus:} Attention weights concentrate on the first token: $W_{\text{attn},S,bsi} \approx \delta_{s0}$ (Kronecker delta)
\item \textbf{Direct Update:} $\mathbf{W}_{\text{update}} = [\mathbf{I}, \mathbf{0}]$ (ignores memory context)
\end{itemize}

Under these conditions:
\begin{align}
\mathbf{A}_{\text{mem}} &\approx \mathbf{x} \quad \text{(identity attention)} \\
\mathbf{s} &\approx \mathbf{x}[:, 0, :] = \mathbf{r} \quad \text{(first token focus)} \\
\mathbf{g}_{\text{attn},\sigma} &\approx \sigma(\mathbf{r} \mathbf{W}_{\text{gate}}) = \mathbf{g}_\sigma \quad \text{(direct gating)}
\end{align}

\subsubsection{Conceptual Interpretation}

\begin{itemize}
\item \textbf{Strategy 1:} "Update memory based on the first token and current memory state"
\item \textbf{Strategy 2:} "Update memory based on attention-weighted input and attention-driven gating"
\end{itemize}

Strategy 2 generalizes Strategy 1 by replacing:
\begin{enumerate}
\item \textbf{Fixed representative} (first token) → \textbf{Attention-weighted representative} (learned selection)
\item \textbf{Direct gating} (input-based) → \textbf{Attention-based gating} (context-aware)
\item \textbf{Simple update signal} (concatenated input+memory) → \textbf{Attention-enhanced signal} (learned combination)
\end{enumerate}

\subsection{Strategy 3 as an Extension of Strategy 2}

\subsubsection{Mathematical Relationship}

Strategy 3 (Sophisticated Update) extends Strategy 2 by adding three key components: memory decay, importance-weighted representation, and adaptive learning rates.

\textbf{Strategy 2 Base:}
\begin{align}
\mathbf{A}_{\text{mem}}, \mathbf{W}_{\text{attn},S} &= \text{MultiHeadAttention}(\mathbf{x}, \mathbf{M}, \mathbf{M}) \\
\mathbf{s} &= \frac{1}{S}\sum_{s=1}^{S} \mathbf{A}_{\text{mem}}[:, s, :] \\
\mathbf{g}_{\text{attn},\sigma} &= \sigma\left(\frac{1}{S}\sum_{s=1}^{S} \mathbf{W}_{\text{attn},S}[:, s, :]\right) \\
\mathbf{M} &\leftarrow (1 - \bar{\mathbf{g}}_{\text{attn},\sigma} \mathbf{1}^T) \odot \mathbf{M} + (\bar{\mathbf{g}}_{\text{attn},\sigma} \mathbf{1}^T) \odot (\mathbf{1} \bar{\mathbf{s}}^T)
\end{align}

\textbf{Strategy 3 Extensions:}
\begin{enumerate}
\item \textbf{Memory Decay:}
\begin{equation}
\mathbf{M} \leftarrow \mathbf{M} \odot (\boldsymbol{\alpha} \mathbf{1}^T)
\end{equation}

\item \textbf{Importance-Weighted Representative:}
\begin{align}
\mathbf{W}_{\text{imp},S} &= \text{softmax}(\sigma(\mathbf{x} \mathbf{W}_{\text{imp}}), \text{dim}=1) \\
\mathbf{r}_{\text{imp}} &= \sum_{s=1}^{S} \mathbf{W}_{\text{imp},S}[:, s] \odot \mathbf{x}[:, s, :]
\end{align}

\item \textbf{Enhanced Memory Interaction:}
\begin{align}
\mathbf{A}_{\text{soph}}, \mathbf{W}_{\text{soph},S} &= \text{MultiHeadAttention}(\mathbf{r}_{\text{imp}}, \mathbf{M}, \mathbf{M})
\end{align}

\item \textbf{Adaptive Learning Rates:}
\begin{align}
\mathbf{g}_{\text{base},\sigma} &= \sigma(\mathbf{r}_{\text{imp}} \mathbf{W}_{\text{gate}}) \\
\mathbf{g}_{\text{final}} &= \mathbf{g}_{\text{base},\sigma} \odot \mathbf{W}_{\text{soph},S}[:, 0, :] \\
\boldsymbol{\lambda} &= \mathbf{g}_{\text{final}} \odot (\mathbf{1} (1 - \boldsymbol{\alpha})^T)
\end{align}
\end{enumerate}

\textbf{Reduction Conditions:}

Strategy 3 reduces to Strategy 2 when:
\begin{itemize}
\item \textbf{No Decay:} $\boldsymbol{\alpha} = \mathbf{1}$ (all ones vector)
\item \textbf{Uniform Importance:} $\mathbf{W}_{\text{imp},S} = \frac{1}{S} \mathbf{1}$ (equal weights for all tokens)
\item \textbf{Identity Memory Interaction:} $\mathbf{W}_{\text{soph},S} = \mathbf{1}$ (uniform attention)
\end{itemize}

Under these conditions:
\begin{align}
\mathbf{r}_{\text{imp}} &= \frac{1}{S}\sum_{s=1}^{S} \mathbf{x}[:, s, :] = \mathbf{s} \quad \text{(uniform averaging)} \\
\mathbf{g}_{\text{final}} &= \mathbf{g}_{\text{base},\sigma} = \sigma(\mathbf{s} \mathbf{W}_{\text{gate}}) = \mathbf{g}_{\text{attn},\sigma} \quad \text{(direct gating)} \\
\boldsymbol{\lambda} &= \mathbf{g}_{\text{final}} \odot \mathbf{0} = \mathbf{0} \quad \text{(no updates due to decay)}
\end{align}

\subsubsection{Conceptual Interpretation}

Strategy 3 adds \textbf{biological realism} and \textbf{temporal awareness}:

\begin{itemize}
\item \textbf{Memory Decay:} Simulates natural forgetting where memories fade over time unless reinforced
\item \textbf{Importance Weighting:} Learns which parts of the input are most relevant for memory updates
\item \textbf{Adaptive Learning Rates:} Combines importance, attention, and decay to determine update strength
\item \textbf{Temporal Dynamics:} Creates a memory system that evolves over time
\end{itemize}

\subsection{Hierarchical Progression}

The three update strategies form a \textbf{hierarchical progression} of increasing sophistication:

\begin{table}[h]
\centering
\begin{tabular}{|l|c|c|c|}
\hline
\textbf{Feature} & \textbf{Strategy 1} & \textbf{Strategy 2} & \textbf{Strategy 3} \\
\hline
Representative Selection & Fixed (first token) & Attention-weighted & Importance-weighted \\
Gating Mechanism & Input-based & Attention-based & Adaptive (multi-factor) \\
Memory Processing & Direct & Direct & Decayed \\
Temporal Awareness & No & No & Yes \\
Context Filtering & Basic & Attention & Importance + Attention \\
Learning Rate & Fixed & Fixed & Adaptive \\
\hline
\end{tabular}
\caption{Hierarchical progression of update strategies}
\end{table}

\subsection{Unified Mathematical Framework}

All three strategies can be expressed in a unified framework:

\begin{align}
\mathbf{M} &\leftarrow \mathbf{M} \odot (\boldsymbol{\alpha} \mathbf{1}^T) \quad &\text{(Memory decay)} \\
\mathbf{r} &= \text{Representative}(\mathbf{x}, \mathbf{M}) \quad &\text{(Representative selection)} \\
\mathbf{g} &= \text{Gating}(\mathbf{r}, \mathbf{M}) \quad &\text{(Update gating)} \\
\mathbf{u} &= \text{UpdateSignal}(\mathbf{r}, \mathbf{M}) \quad &\text{(Update signal)} \\
\mathbf{M} &\leftarrow (1 - \bar{\mathbf{g}} \mathbf{1}^T) \odot \mathbf{M} + (\bar{\mathbf{g}} \mathbf{1}^T) \odot (\mathbf{1} \bar{\mathbf{u}}^T) \quad &\text{(Memory update)}
\end{align}

Where:
\begin{itemize}
\item \textbf{Strategy 1:} $\boldsymbol{\alpha} = \mathbf{1}$, Representative = first token, Gating = direct, UpdateSignal = concatenated
\item \textbf{Strategy 2:} $\boldsymbol{\alpha} = \mathbf{1}$, Representative = attention-weighted, Gating = attention-based, UpdateSignal = attention-enhanced
\item \textbf{Strategy 3:} $\boldsymbol{\alpha}$ = learned, Representative = importance-weighted, Gating = adaptive, UpdateSignal = context-aware
\end{itemize}

\subsection{Memory Storage Dynamics}

\subsubsection{Strategy 1: Static Storage}
\begin{itemize}
\item \textbf{Storage Pattern:} Fixed updates based on first token
\item \textbf{Memory Evolution:} Linear combination of old and new information
\item \textbf{Forgetting:} No explicit forgetting mechanism
\item \textbf{Context Sensitivity:} Limited to current input and average memory state
\end{itemize}

\subsubsection{Strategy 2: Attention-Driven Storage}
\begin{itemize}
\item \textbf{Storage Pattern:} Context-aware updates based on attention patterns
\item \textbf{Memory Evolution:} Attention-weighted combination of old and new information
\item \textbf{Forgetting:} No explicit forgetting, but attention can suppress updates
\item \textbf{Context Sensitivity:} High - uses full sequence context for updates
\end{itemize}

\subsubsection{Strategy 3: Adaptive Storage}
\begin{itemize}
\item \textbf{Storage Pattern:} Biologically inspired updates with temporal dynamics
\item \textbf{Memory Evolution:} Decay + importance + attention + adaptive learning rates
\item \textbf{Forgetting:} Explicit decay mechanism with learnable rates
\item \textbf{Context Sensitivity:} Maximum - uses importance weighting + attention + temporal awareness
\end{itemize}

\subsection{Design Implications}

This hierarchical relationship has important implications for memory system design:

\begin{enumerate}
\item \textbf{Progressive Complexity:} Each strategy builds upon the previous one, allowing for gradual sophistication
\item \textbf{Memory Efficiency:} Simpler strategies use less computational resources
\item \textbf{Biological Plausibility:} More sophisticated strategies better model human memory processes
\item \textbf{Task Adaptability:} Different tasks may require different levels of memory sophistication
\item \textbf{Training Stability:} Simpler strategies may be easier to train initially
\end{enumerate}

\subsection{Memory Storage Requirements}

All three strategies require \textbf{persistent memory storage} during training:

\begin{itemize}
\item \textbf{Strategy 1:} Memory matrix $\mathbf{M}$ must persist across training steps
\item \textbf{Strategy 2:} Memory matrix + attention parameters must persist
\item \textbf{Strategy 3:} Memory matrix + decay factors + attention parameters must persist
\end{itemize}

Without persistent storage, there would be no memory to retrieve or update, making the entire memory mechanism non-functional.

The unified framework demonstrates that memory update sophistication lies on a continuum, where each strategy represents a different point balancing computational cost, biological plausibility, and representational power for memory storage and consolidation.

\section{Appendix E: Temporal Memory Decay and Relative Position Dynamics}

This appendix explores the reformulation of memory decay mechanisms from static decay factors to dynamic, position-dependent decay functions that better model biological memory processes and enable more sophisticated temporal reasoning in neural architectures.

\subsection{From Static to Relative Decay}

Traditional memory decay mechanisms employ static decay factors $\alpha_i \in [0,1]$ that remain constant throughout the sequence. While computationally efficient, this approach fails to capture the fundamental temporal dynamics of human memory, where forgetting rates depend on the temporal distance between memory formation and retrieval. Recent advances in positional encoding research have demonstrated that relative position information can significantly enhance long-context modeling capabilities.

The reformulation introduces **relative position decay** where memory strength depends on the temporal distance between memory formation and current processing. Let $t$ represent the current sequence position and $t_i$ denote the position where memory slot $i$ was last updated. The decay function becomes $\alpha(t, t_i) = \alpha(|t - t_i|)$, where the decay rate depends on the relative temporal distance $\Delta t = |t - t_i|$.

This temporal formulation enables several biologically plausible decay patterns. **Exponential decay** models the rapid initial forgetting observed in human memory: $\alpha(\Delta t) = \exp(-\lambda \Delta t)$, where $\lambda$ controls the decay rate. **Power law decay** captures the long-tail forgetting patterns characteristic of semantic memory: $\alpha(\Delta t) = (1 + \Delta t)^{-\beta}$, where $\beta$ determines the decay exponent. **Sigmoid decay** provides a more gradual transition between short-term and long-term memory: $\alpha(\Delta t) = 1 / (1 + \exp(\gamma(\Delta t - \tau)))$, where $\gamma$ controls the transition steepness and $\tau$ determines the transition point.

\subsection{Integration with Positional Encoding Research}

Recent work in positional encoding has revealed important insights about the relationship between position and memory in transformer architectures. The **3D Rotary Position Encoding (3D-RPE)** \cite{3d-rpe2024} extends traditional 2D RoPE by mapping positional encodings onto a three-dimensional sphere, enabling controllable long-term decay within specified chunk sizes. This approach demonstrates that explicit decay mechanisms can be integrated into attention-based architectures while maintaining computational efficiency.

The **High-frequency rotary Position Encoding (HoPE)** \cite{hope2024} challenges the traditional long-term decay assumption in positional encodings by replacing specific components with position-independent ones. While HoPE removes decay constraints to enhance context awareness, it highlights the importance of understanding how temporal dynamics affect memory retrieval in long sequences.

**Retentive Networks** \cite{retnet2023} incorporate decay directly into the attention mechanism through exponential decay masks, providing a different approach to temporal memory modeling. The attention mask in Retentive Networks applies exponential decay based on relative position: $A_{ij} = \exp(-\gamma |i - j|)$, where $\gamma$ controls the decay rate. This approach demonstrates that decay mechanisms can be effectively integrated into attention-based architectures without requiring separate memory storage.

\subsection{Physical Time Integration}

The integration of physical wall-clock time introduces additional temporal dynamics that can enhance memory systems. By associating memory updates with physical timestamps, the system can model memory decay based on real-world temporal patterns rather than sequence position alone. This approach enables the modeling of memory consolidation processes that occur over extended periods, independent of sequence processing.

Let $T_{\text{update}}^{(i)}$ represent the wall-clock time when memory slot $i$ was last updated, and $T_{\text{current}}$ denote the current processing time. The physical time decay function becomes $\alpha_{\text{physical}}(T_{\text{current}}, T_{\text{update}}^{(i)}) = \alpha(T_{\text{current}} - T_{\text{update}}^{(i)})$, where the decay depends on the actual temporal distance between updates.

This physical time integration enables several sophisticated memory behaviors. **Circadian decay patterns** can model the influence of sleep-wake cycles on memory consolidation, where memories formed during different phases of the circadian rhythm exhibit different decay characteristics. **Context-dependent decay** can incorporate environmental factors such as stress levels, attention state, or emotional context that influence memory retention. **Multi-scale temporal dynamics** can simultaneously model both rapid forgetting (seconds to minutes) and long-term consolidation (hours to days) within the same memory system.

\subsection{Hybrid Temporal Decay Models}

The combination of sequence position and physical time enables hybrid decay models that capture both immediate contextual effects and long-term consolidation processes. The **dual-time decay function** combines both temporal scales: $\alpha_{\text{hybrid}}(t, T) = \alpha_{\text{sequence}}(t) \cdot \alpha_{\text{physical}}(T)$, where sequence-based decay captures immediate contextual forgetting and physical time decay models long-term consolidation.

This hybrid approach allows the memory system to distinguish between **working memory decay** (rapid forgetting over sequence positions) and **episodic memory consolidation** (slower decay over physical time). The system can learn different decay parameters for different types of information, enabling sophisticated memory management that adapts to the temporal characteristics of the stored information.

\subsection{Implementation Considerations}

The implementation of relative position decay requires careful consideration of computational efficiency and memory management. **Precomputed decay tables** can store decay values for common temporal distances, reducing computational overhead during inference. **Adaptive decay scheduling** can adjust decay rates based on the current context or task requirements, enabling dynamic memory management.

**Memory consolidation mechanisms** can periodically update decay parameters based on the importance and recency of stored information. This adaptive approach ensures that critical memories receive appropriate protection from decay while allowing less important information to fade naturally. The system can learn optimal decay schedules through gradient-based optimization, enabling task-specific memory management strategies.

\subsection{Biological Plausibility and Future Directions}

The relative position decay formulation provides a more biologically plausible model of memory processes than static decay factors. Human memory exhibits complex temporal dynamics that depend on both recency and importance, with different types of information following different forgetting curves. The proposed framework enables the modeling of these sophisticated temporal patterns within neural architectures.

Future research directions include the investigation of **hierarchical temporal decay** where different memory layers exhibit different temporal characteristics, **context-dependent decay modulation** where external factors influence forgetting rates, and **adaptive consolidation** where the system learns optimal decay schedules for different types of information. These extensions would further enhance the biological plausibility and practical utility of temporal memory systems in neural architectures.

The integration of relative position decay with existing positional encoding techniques and physical time dynamics represents a significant advancement in the development of biologically inspired memory systems for neural architectures. This approach provides a foundation for more sophisticated temporal reasoning capabilities that can enhance the performance of large language models and other sequence processing systems.

\section{Appendix F: Importance Weighting in Memory Retrieval}

This appendix provides a detailed analysis of the importance weighting mechanism employed in Strategy 3 (Sophisticated Retrieval), explaining its computational implementation, biological motivation, and practical benefits for context-aware memory systems.

\subsection{Mechanism Overview}

The importance weighting mechanism serves as a **final contextual filter** that dynamically adjusts the strength of retrieved memory based on the current input's characteristics. Unlike the attention mechanisms that determine which memory slots to retrieve from, importance weighting determines how much of the retrieved information should actually be used in the current context.

The mechanism operates through a two-step process. First, the system computes importance scores for each input token using a learned projection matrix. Second, these scores are applied as multiplicative gates to the retrieved memory representations, effectively filtering the retrieved information based on its perceived relevance to the current context.

\subsection{Mathematical Formulation}

The importance weighting mechanism is implemented through the following mathematical operations:

\begin{align}
\mathbf{I}_{\text{weights},\sigma} &= \sigma(\mathbf{x} \mathbf{W}_{\text{imp}}) \in \mathbb{R}^{B \times S \times 1} \\
I_{\text{weights},\sigma,bs1} &= \sigma\left(\sum_{j=1}^{D} x_{bsj} W_{\text{imp},j1}\right) \\
\mathbf{R}_{\text{mem}} &= \mathbf{R}_{\text{mem}} \odot \mathbf{I}_{\text{weights},\sigma} \in \mathbb{R}^{B \times S \times D} \\
R_{\text{mem},bsj} &= R_{\text{mem},bsj} I_{\text{weights},\sigma,bs1}
\end{align}

The first equation computes importance scores by projecting the input sequence through a learned matrix $\mathbf{W}_{\text{imp}} \in \mathbb{R}^{D \times 1}$ and applying a sigmoid activation. The sigmoid ensures that importance scores are bounded between 0 and 1, where values closer to 1 indicate high relevance and values closer to 0 indicate low relevance. The second equation applies these importance scores as multiplicative gates to the retrieved memory, with broadcasting ensuring that each importance score affects all dimensions of the corresponding retrieved memory vector.

\subsection{Biological Motivation}

The importance weighting mechanism is inspired by **selective attention** processes in human cognition. When humans retrieve information from memory, they do not use all retrieved information equally. Instead, they selectively focus on information that is most relevant to the current task or context, effectively filtering out irrelevant details while amplifying important information.

This selective attention process is crucial for efficient cognitive processing. Without it, humans would be overwhelmed by irrelevant information retrieved from memory, making it difficult to focus on the task at hand. The importance weighting mechanism provides a computational analog of this biological process, enabling neural networks to exhibit similar selective attention behaviors.

The sigmoid activation function used in importance weighting is particularly well-suited for this purpose. It provides a smooth, differentiable way to implement binary-like decisions about information relevance while maintaining gradient flow for learning. The bounded output range ensures that the mechanism can completely suppress irrelevant information (scores near 0) or fully utilize important information (scores near 1), with smooth transitions between these extremes.

\subsection{Computational Benefits}

The importance weighting mechanism provides several computational advantages that enhance the overall performance of memory-augmented neural networks. **Contextual relevance filtering** enables the system to dynamically adjust its use of retrieved information based on the current input, preventing the system from being distracted by irrelevant memories. This filtering is particularly important in long sequences where the system may retrieve information that was relevant earlier in the sequence but is no longer useful for the current processing step.

**Dynamic memory strength modulation** allows the system to amplify important memories while suppressing less relevant ones. This capability is crucial for maintaining focus on the most important information, especially when the memory system retrieves multiple pieces of information that may have varying degrees of relevance to the current context. The system can learn to identify which types of information are most useful for different types of inputs, enabling sophisticated context-dependent memory usage.

**Computational efficiency** is another key benefit of importance weighting. By reducing the effective strength of less relevant retrieved information, the mechanism can reduce the computational impact of irrelevant memories on downstream processing. This efficiency gain is particularly important in large-scale systems where memory retrieval may return substantial amounts of information that may not all be equally useful.

\subsection{Practical Applications}

The importance weighting mechanism has broad applicability across various domains where context-aware memory systems are beneficial. In **natural language processing**, the mechanism can help systems focus on the most relevant retrieved information when processing complex texts, improving the quality of generated responses and reducing the influence of irrelevant retrieved context.

In **question answering systems**, importance weighting can help the system focus on the most relevant retrieved passages when answering questions, improving accuracy and reducing the impact of irrelevant retrieved information. The mechanism enables the system to dynamically adjust its attention to different parts of the retrieved information based on the specific question being asked.

**Conversational AI systems** can benefit from importance weighting by enabling more contextually appropriate responses. The mechanism allows the system to focus on the most relevant aspects of the conversation history when generating responses, improving coherence and relevance while reducing the influence of less relevant historical information.

\subsection{Learning Dynamics}

The importance weighting mechanism learns to identify relevant information through gradient-based optimization. The learned projection matrix $\mathbf{W}_{\text{imp}}$ develops representations that can distinguish between relevant and irrelevant retrieved information based on the current input context. This learning process enables the system to develop sophisticated understanding of which types of information are most useful in different contexts.

The sigmoid activation function provides smooth gradients that facilitate learning, enabling the system to gradually refine its understanding of information relevance. The bounded output range ensures that the learning process remains stable while allowing the system to develop strong preferences for relevant information and strong suppression of irrelevant information.

\subsection{Integration with Other Memory Mechanisms}

The importance weighting mechanism integrates seamlessly with other memory mechanisms in the sophisticated retrieval strategy. It operates after memory decay and attention-based retrieval, providing a final layer of contextual filtering that enhances the overall effectiveness of the memory system. The mechanism complements the attention mechanisms by providing additional context-dependent filtering that operates at a different level of the memory processing pipeline.

The integration with RMS normalization ensures that the importance weighting mechanism operates on properly normalized retrieved memory representations, preventing numerical instabilities that could arise from the multiplicative gating operation. This integration creates a robust memory processing pipeline that can handle complex contextual filtering while maintaining computational stability.

The importance weighting mechanism represents a crucial component of sophisticated memory systems, providing the contextual awareness necessary for effective memory usage in complex neural architectures. Its biological motivation, computational benefits, and practical applications make it an essential tool for developing more intelligent and contextually aware artificial intelligence systems.

\section{Appendix G: Normalization Strategies in Attention and Memory Mechanisms}

This appendix examines the choice between softmax and sigmoid normalization in attention mechanisms and memory systems, exploring the theoretical foundations, biological motivations, and practical implications of different normalization strategies.

\subsection{The Softmax-Sigmoid Dichotomy}

The choice between softmax and sigmoid normalization in neural architectures has traditionally been based on different use cases and computational requirements. Softmax normalization, defined as $\text{softmax}(x_i) = \exp(x_i) / \sum_j \exp(x_j)$, creates a probability distribution where weights sum to unity. Sigmoid normalization, defined as $\sigma(x_i) = 1 / (1 + \exp(-x_i))$, produces independent bounded values between 0 and 1. While both approaches have their merits, the biological and computational considerations suggest that sigmoid normalization may be more appropriate for many memory and attention applications.

\subsection{Biological Motivation for Sigmoid Normalization}

From a biological perspective, sigmoid normalization better reflects the independent nature of attention and memory processes in human cognition. When humans attend to multiple stimuli or retrieve multiple memories, the strength of attention to one stimulus does not necessarily reduce the attention paid to others. This independence is naturally captured by sigmoid normalization, where each attention weight is computed independently and can take any value between 0 and 1.

The sigmoid function itself has strong biological roots, as it closely approximates the activation function of biological neurons. The smooth, differentiable nature of the sigmoid function enables gradient-based learning while maintaining the bounded output range that prevents numerical instabilities. This biological plausibility makes sigmoid normalization particularly suitable for memory systems that aim to model human cognitive processes.

\subsection{Computational Efficiency Considerations}

Sigmoid normalization offers significant computational advantages over softmax normalization. The computation of sigmoid values requires only a single exponential operation per element, while softmax requires computing exponentials for all elements followed by a normalization step. This difference becomes particularly important in large-scale systems where attention mechanisms operate over many memory slots or sequence positions.

The independence of sigmoid values also enables more efficient parallel computation, as each attention weight can be computed independently without requiring knowledge of other weights. This parallelizability is crucial for GPU acceleration and can significantly improve training and inference speed in large neural networks.

\subsection{Alternative Normalization Strategies}

Beyond the traditional softmax-sigmoid dichotomy, several alternative normalization strategies can provide different trade-offs between computational efficiency and representational power. **Average normalization** divides the sum of sigmoid values by the number of elements, providing bounded outputs while maintaining independence: $\text{avg\_norm}(x_i) = \sigma(x_i) / \frac{1}{N}\sum_{j=1}^N \sigma(x_j)$.

**Variance-stabilized normalization** divides by the square root of the number of elements to maintain bounded variance, similar to the scaling used in attention mechanisms: $\text{var\_norm}(x_i) = \sigma(x_i) / \sqrt{N}$. This approach ensures that the variance of the normalized values remains constant regardless of the sequence length, preventing the attention weights from becoming too small in long sequences.

**Learnable normalization** introduces a learnable parameter that can adapt the normalization strength during training: $\text{learnable\_norm}(x_i) = \sigma(x_i) / \alpha$, where $\alpha$ is a learnable parameter. This approach allows the model to learn the optimal normalization strategy for its specific task and data distribution.

\subsection{Attention Mechanism Normalization}

In traditional attention mechanisms, softmax normalization serves to create a probability distribution over the attention targets. However, this normalization constraint may not always be necessary or beneficial. The exponential function in the numerator of softmax can be replaced by any positive, bounded function without fundamentally changing the attention mechanism's behavior.

**Sigmoid-based attention** replaces the exponential function with sigmoid: $A_{ij} = \sigma(s_{ij}) / \sum_k \sigma(s_{ik})$, where $s_{ij}$ are the raw attention scores. This approach maintains the probability distribution property while using the more biologically plausible sigmoid function.

**Independent attention weights** eliminate the normalization constraint entirely, using sigmoid values directly: $A_{ij} = \sigma(s_{ij})$. This approach allows for more flexible attention patterns where the total attention strength can vary based on the input context, potentially enabling more sophisticated attention behaviors.

\subsection{Memory System Normalization}

In memory systems, the choice of normalization strategy can significantly impact the system's ability to learn and adapt. **Importance weighting** in memory retrieval can benefit from sigmoid normalization, as it allows for independent assessment of each memory slot's relevance without requiring the weights to sum to unity.

**Memory update gating** can also benefit from sigmoid normalization, as it enables independent control over each memory slot's update strength. This independence is particularly important in sophisticated memory systems where different memory slots may require different update strategies based on their content and relevance.

\subsection{Practical Implementation Considerations}

The implementation of alternative normalization strategies requires careful consideration of numerical stability and computational efficiency. **Numerical stability** can be maintained by using the log-sum-exp trick for softmax alternatives or by ensuring that sigmoid values are computed with sufficient precision.

**Gradient flow** must be preserved to enable effective learning. Sigmoid normalization provides smooth gradients that facilitate learning, while alternative normalization strategies may require additional care to ensure that gradients flow properly through the normalization operation.

**Memory efficiency** is another important consideration, particularly in large-scale systems. Sigmoid normalization requires less memory than softmax normalization, as it does not require storing intermediate normalization values. This memory efficiency can be crucial in systems with limited memory resources.

\subsection{Experimental Validation}

The choice of normalization strategy should be validated through empirical experimentation. **Ablation studies** can compare different normalization strategies on the same task to determine which approach provides the best performance for the specific application.

**Computational profiling** can measure the actual performance differences between normalization strategies, taking into account both computational cost and memory usage. This profiling should be performed on the target hardware to ensure that the theoretical advantages translate to practical benefits.

**Biological plausibility testing** can evaluate whether sigmoid-based normalization produces more human-like attention and memory patterns. This evaluation can be performed through behavioral experiments or by comparing the learned representations to those observed in biological systems.

\subsection{Future Directions}

The exploration of alternative normalization strategies opens several promising research directions. **Adaptive normalization** could allow the system to learn the optimal normalization strategy for different contexts or tasks. **Hierarchical normalization** could apply different normalization strategies at different levels of the memory hierarchy.

**Biological validation** could provide insights into which normalization strategies best match human cognitive processes, potentially leading to more biologically plausible artificial intelligence systems. **Theoretical analysis** could provide deeper understanding of the mathematical properties of different normalization strategies and their implications for learning and generalization.

The choice of normalization strategy represents a fundamental design decision that can significantly impact the performance and biological plausibility of memory and attention systems. While traditional approaches have relied heavily on softmax normalization, the biological and computational advantages of sigmoid normalization suggest that alternative approaches may provide better solutions for many applications.

\end{document}
